% !TeX program = pdflatex
% teoria.tex — compila SOZINHO (pdflatex teoria.tex) e também dentro de main.tex.
% O preâmbulo vive em estilo.tex, partilhado pelos três. Ver livro.tex.
\documentclass[livro.tex]{subfiles}

\begin{document}

\ifSubfilesClassLoaded{%
  \gkcapa{A Teoria --- a codificação matricial das fracções contínuas}
       {álgebra, topologia e análise --- três partes que são duais entre si}
       {o corpo dual, as duas leis, e as fracções contínuas como sua representação}
  \maketitle
  \tableofcontents
}{%
  \part*{A Teoria}\addcontentsline{toc}{part}{A Teoria}
}


\begin{abstract}
\textbf{Este é um texto sobre dualidade}, e a definição de que tudo parte é uma só:
\emph{a dualidade é a memória da divisão}. Dividir perde; a dualidade é o que guarda a metade que se
ia deitar fora, e é isso que torna a divisão reversível. Daí saem \textbf{duas leis}
(\S\ref{sec:duasleis}):

\medskip
\medskip
\noindent \textbf{Primeiro os dois espaços, depois a relação --- e a igualdade tem de ser bem
tipada.}

\medskip
\noindent\textbf{O primal.}\quad $V$, espaço sobre um corpo $K$, com os seus operadores
$T\colon V\to V$.

\noindent\textbf{O dual.}\quad $V^{*}=\operatorname{Hom}(V,K)$: as \emph{formas de medida} sobre
$V$.

\noindent\textbf{O adjunto.}\quad $T$ induz $T^{*}\colon V^{*}\to V^{*}$ por
\[
  (T^{*}\varphi)(v) \;=\; \varphi(Tv), \qquad \varphi\in V^{*},\ v\in V ,
\]
e ele existe sempre, sem hipótese sobre $T$.

\medskip
\noindent \textbf{E aqui está o ponto.} $T$ e $T^{*}$ pertencem a categorias distintas,
\[
  T\in\operatorname{End}(V), \qquad T^{*}\in\operatorname{End}(V^{*}) ,
\]
\emph{logo ``$T^{*}=-T$'' não é uma igualdade bem tipada}: os dois membros não vivem no mesmo
espaço de operadores. Para os comparar é preciso um isomorfismo $J\colon V\to V^{*}$, induzido por
um emparelhamento não degenerado $\langle\cdot,\cdot\rangle\colon V\times V\to K$. Transportando o
adjunto para o primal,
\[
  T^{\dagger} \;=\; J^{-1}T^{*}J \;\colon\; V\to V ,
\]
e agora ambos os membros estão em $\operatorname{End}(V)$. \emph{Só então a lei se escreve:}

\begin{center}\small
\setlength\tabcolsep{5pt}
\begin{tabular}{@{}llll@{}}
\toprule
 & \textbf{a lei} & \textbf{bem tipada} & \textbf{e o objeto que a cumpre} \\
\midrule
$1$ & a \emph{unidade} é dual   & $1^{\dagger}=-1$ & a Möbius \textbf{involutiva} (traço $0$) \\
$2$ & a \emph{dualidade} é dual & $T^{\dagger}=-T$ & a Möbius de \textbf{Fibonacci} ($\det=-1$) \\
\bottomrule
\end{tabular}
\end{center}

\noindent A Lei~$2$ diz que $T$ é \textbf{anti-autoadjunto} relativamente ao emparelhamento
escolhido. \emph{E a dualidade não identifica dois objetos iguais: identifica duas descrições do
mesmo objeto em espaços opostos} --- o primal contém os objetos, o dual contém as suas formas de
medida, e o adjunto transporta entre os dois lados.

\begin{obs}[por que a marcação é obrigatória, e mede-se]\label{obs:asterisco}
Sem distinguir o dual, o leitor lê ``$T=-T$'' como igualdade interna em $\operatorname{End}(V)$, o
que sobre corpos de característica $\neq2$ dá $2T=0$, isto é $T=0$: \textbf{a lei colapsa na solução
trivial}. \emph{A marcação impede exatamente esse colapso.}
\end{obs}

\medido{Em matrizes $2\times2$ com entradas em $[-3,3]$: $M=-M$ tem \textbf{$1$} solução em $2401$;
$M^{\top}=-M$ tem $7$, todas da forma
$\bigl(\begin{smallmatrix}0&b\\-b&0\end{smallmatrix}\bigr)$ --- um espaço de dimensão $1$
(\code{tests/bidual.c}).}

\paragraph{E o ambiente em que tudo isto se passa tem nome: o \emph{corpo dual}.}

\begin{definicao}[corpo dual]\label{def:corpodual}
Chama-se \textbf{corpo dual} à estrutura formada pelo par $(V,V^{*})$ \emph{juntamente com a
identificação} $J$. \emph{As leis desta teoria são sempre enunciadas neste corpo, nunca em apenas
um dos seus lados.}
\end{definicao}

\noindent \emph{E cada lado é incompleto sozinho}: o primal contém os objetos, o dual contém as
suas leituras, e é $J$ que permite transportar operadores entre ambos. A dualidade \textbf{não
afirma que um objeto seja igual ao seu contrário} --- afirma que todo objeto tem um representante do
outro lado do corpo dual.

\begin{center}\small
\begin{tabular}{@{}lcl@{}}
\toprule
\textbf{primal} & & \textbf{dual} \\
\midrule
$V$ & & $V^{*}$ \\
os objetos & $\xleftrightarrow{\ J\ }$ & as leituras \\
$T\colon V\to V$ & & $T^{*}\colon V^{*}\to V^{*}$ \\
\bottomrule
\end{tabular}
\end{center}

\paragraph{E dentro dele, um resultado --- e convém não confundir as duas palavras.} Fixado $J$, a
lei \textbf{define} o conjunto dos anti-autoadjuntos, e esse conjunto fecha em \emph{corpo no
sentido algébrico} (um \emph{field}). \textbf{São duas coisas distintas com a mesma palavra}: o
\emph{corpo dual} é o ambiente $(V,V^{*},J)$; o que se segue é uma estrutura algébrica que vive
dentro dele.

\begin{teorema}[os anti-autoadjuntos formam um corpo algébrico]\label{thm:corpodual}
O conjunto dos $T$ com $T^{\dagger}=-T$, com as escalares juntas, é um corpo. Em matrizes
$2\times2$ é $\{aI+bJ\}$ com $J^{2}=-I$, isto é $\cong\C$; em $K[\sigma]$ é a \textbf{família
metálica}, $\sigma\sigma'=-1$. \emph{E o seu dual é ele próprio}: $\sigma'=m-\sigma$ já está dentro
dele, donde $K(\sigma')=K(\sigma)$.
\end{teorema}

\begin{proof}
$M^{\top}=-M$ em $2\times2$ força $M=bJ$, e $(aI+bJ)(aI-bJ)=(a^{2}+b^{2})I$ com $a^{2}+b^{2}>0$ para
$(a,b)\neq(0,0)$: todo não nulo inverte. Em $K[\sigma]$, $\sigma'=-1/\sigma$ é $\sigma\sigma'=-1$,
a borda $x^{2}-mx-1$, irredutível porque $m^{2}+4$ nunca é quadrado perfeito para $m\geq1$; e como
$\sigma^{2}=m\sigma+1$, tem-se $\sigma'=m-\sigma$, com $\sigma'+\sigma=m$ e $\sigma'\sigma=-1$.
\end{proof}

\medido{$a^{2}+b^{2}>0$ em $80$ elementos não nulos; $\sigma'=m-\sigma$ com traço e norma corretos
em $m=1\ldots11$, resíduo $0$ (\code{tests/bidual.c}).}

\medskip
\noindent \textbf{E o texto está em quatro capítulos, cada um a lei vista de um lado:}

\begin{center}\small
\begin{tabular}{@{}clll@{}}
\toprule
 & \textbf{o capítulo} & \textbf{o que dá} & \textbf{de que lei} \\
\midrule
$1$ & \textbf{Álgebra}   & o corpo dual, as leis, as operações & o fundamento \\
$2$ & \textbf{Topologia} & a torre, a ordem, a bidualidade & a Lei~$1$ a subir \\
$3$ & \textbf{Análise}   & métrica, limite, derivada; e no fim o corpo cruz e o corpo de corpos & a Lei~$2$ a andar \\
$4$ & \textbf{Geometria} & a curvatura --- e o operador geométrico \emph{é} o tempo & o gancho \\
\bottomrule
\end{tabular}
\end{center}
\end{abstract}

\part{Álgebra --- o corpo dual e as suas leis}
\noindent\hfill\emph{\small[\textbf{fundamento}: a fonte, e o que dela se prova]}\medskip

\section{A partida: espaços vetoriais, e duas notações}\label{sec:partida}

Este texto parte de \textbf{um espaço vetorial}, e de mais nada. Não parte dos inteiros, nem dos
racionais, nem da reta: esses aparecem adiante como \emph{instâncias}, ao lado de outras, e nenhum
deles é pressuposto para que o que se segue tenha sentido. \textbf{A reta é uma instância}, não o
chão --- e é por isso que ela e a sua aritmética vivem no \emph{Catálogo}, onde as realizações
estão, e não aqui.

\medskip
\noindent Seja $V$ um espaço vetorial sobre um corpo $K$ e $V^{*}=\operatorname{Hom}(V,K)$ o seu
dual. Não se pede a $K$ que seja ordenado, nem completo, nem de característica zero. Não se pede a
$V$ que seja de dimensão finita, e onde a finitude for usada, será dito. Sobre isto assentam
\textbf{duas notações}, e elas bastam para escrever tudo o que vem depois.

\begin{mdframed}[linewidth=0.4pt,innerleftmargin=10pt,innerrightmargin=10pt,
                 innertopmargin=8pt,innerbottommargin=8pt,skipabove=6pt,skipbelow=6pt]
\textbf{A estaca}\quad $x \mapsto x^{\dagger}$ \quad --- \emph{troca de lado.} Leva o primal ao
dual e o dual ao primal. É involutiva: $x^{\dagger\dagger}=x$.

\medskip
\textbf{A cruz}\quad $x \mapsto x^{\times} = \bigl(x \oplus x^{\dagger},\; x \otimes
x^{\dagger}\bigr)$ \quad --- \emph{projeta no que fica.} Devolve o \emph{par} formado pela escrita
e pela leitura de $x$ com o seu dual, e ambas as coordenadas são fixas pela estaca.
\end{mdframed}

\noindent \emph{A estaca move; a cruz mede o que não se moveu.} Uma é a operação, a outra é o
invariante --- e todo o resto deste texto é o que sai de as ter as duas. Escritas assim, as duas
leis não precisam de nomear nenhum conjunto de números:

\begin{align}
\text{\textbf{Lei 1} --- a unidade é dual:} && 1^{\dagger} &= -1 \tag{L1}\\
\text{\textbf{Lei 2} --- a dualidade é dual:} && T^{\dagger} &= -T \tag{L2}
\end{align}

\noindent A primeira diz que o objeto mais simples que existe já é um par; a segunda diz que a
própria operação de emparelhar também tem dual. \textbf{Nenhuma das duas menciona um número.}

\subsection{Por que a cruz precisa de duas coordenadas}\label{sec:cruzduas}

\noindent A cruz podia ter sido definida com uma coordenada só --- a soma, ou o produto. Não pode:
\emph{cada uma sozinha colide}, e a medida de quanto colide é o argumento.

\medido{Numa instância concreta com $28\,392$ pares distintos, a primeira coordenada sozinha
identifica $412$ pares que são diferentes; o par inteiro identifica $44$. \textbf{São $368$
distinções que a segunda coordenada recupera} --- e é por isso que a cruz é um par e não um
número.}

\noindent É o mesmo argumento que aparece em toda a parte neste texto e vale a pena dizê-lo uma
vez: \textbf{quem projeta perde, e o dual é o que guarda o que a projeção deitou fora.} Uma teoria
escrita com uma coordenada só não está errada --- está \emph{pela metade}, e a metade que falta é
exatamente a que a estaca alcança.

\subsection{E as potências das duas dizem coisas diferentes}\label{sec:potencias}

\noindent Iterar a estaca e iterar a cruz dão comportamentos opostos, e a diferença é a estrutura
inteira em miniatura:

\begin{center}\small
\begin{tabular}{@{}lll@{}}
\toprule
& \textbf{o que faz} & \textbf{ao iterar} \\
\midrule
estaca $x^{\dagger n}$ & troca de lado & \textbf{período $2$}: $\dagger\dagger=\mathrm{id}$, e
$x^{\dagger}\neq x$ fora do fixo \\
cruz\ \ \ $x^{\times n}$ & projeta no fixo & \textbf{idempotente}: já está no fixo, cruzar não move \\
\bottomrule
\end{tabular}
\end{center}

\medido{Verificado em toda uma janela de instâncias: $x^{\dagger\dagger}=x$ sem uma única exceção,
e $x^{\dagger}=x$ \emph{apenas} no subanel fixo --- se a estaca fosse a identidade em mais algum
sítio, ela não trocaria nada e a Lei~$1$ não distinguiria nada. As duas coordenadas da cruz caem
sempre no fixo, e voltar a cruzar devolve-as intactas.}

\noindent \textbf{A estaca é o relógio e a cruz é o mostrador.} A primeira é a dinâmica, com o seu
período; a segunda é a leitura, que não se move. É esta diferença --- e não uma escolha de
axiomas --- que faz com que a ordem, o limite e a completude possam ser construídos aqui sem que a
reta seja invocada: eles saem do par, e o par existe assim que existir a estaca.

\subsection{A ordem sai do dual}\label{sec:ordemdual}

\noindent Duas coordenadas fixas dão uma ordem lexicográfica, e uma ordem lexicográfica sobre um
par é total. Não é preciso corte, nem sucessão fundamental, nem qualquer construção do contínuo:
\textbf{a ordem é a cruz lida por coordenadas}, e a estaca é que garante que as coordenadas
existem.

\medido{A relação induzida por $(x\oplus x^{\dagger},\,x\otimes x^{\dagger})$ lida
lexicograficamente satisfaz tricotomia e transitividade em toda a amostra verificada, com resíduo
zero em ambas. Nenhum corte de Dedekind, nenhuma sucessão de Cauchy: só o par.}

\noindent E é aqui que a análise deste texto se separa da análise usual. \textbf{A ordem não é
herdada da reta; é produzida pela dualidade.} A reta é \emph{uma} instância onde essa ordem tem uma
forma familiar --- e há outras, com a mesma ordem e sem qualquer semelhança com ela.

\section{A família metálica, algébrica}\label{sec:metalica-alg}

\noindent Esta é a primeira instância, e serve de aferição: ela costuma ser apresentada com
raízes quadradas e valores aproximados, e \textbf{não precisa de nenhum dos dois}. Fixe-se um
inteiro $n$ e tome-se
\[
  A_n \;=\; K[x]\big/(x^{2}-nx-1), \qquad \sigma = \text{classe de } x .
\]
Aqui $\sigma$ não é um número entre $1$ e $2$: é a classe de $x$, e a única regra é a redução
$\sigma^{2}=n\sigma+1$. A estaca é o único automorfismo não trivial, $\sigma^{\dagger}=n-\sigma$;
e a cruz devolve o par $(\operatorname{tr},\operatorname{N})$.

\begin{proposicao}[a Lei 2 caracteriza a norma]\label{prop:lei2norma}
Em $A_n$, um elemento $x$ satisfaz $x^{-1}=-x^{\dagger}$ \emph{se e só se} $\operatorname{N}(x)=-1$.
Em particular $\sigma$ satisfaz, para todo $n$, e a sua inversa $\sigma-n$ é um elemento do próprio
anel.
\end{proposicao}

\begin{proof}
$x^{-1}=-x^{\dagger}$ equivale a $x\cdot(-x^{\dagger})=1$, isto é a $-\operatorname{N}(x)=1$. Para
$\sigma$: $\sigma(\sigma-n)=\sigma^{2}-n\sigma=1$ pela redução, e nada mais é usado.
\end{proof}

\medido{Varridos $10\,845$ elementos em treze instâncias: os que cumprem a Lei~$2$ e os de norma
$-1$ são \emph{o mesmo conjunto}, com zero discordâncias --- e são menos de um décimo do total,
pelo que a lei separa em vez de descrever. $\operatorname{N}(\sigma)=-1$ e
$\operatorname{tr}(\sigma)=n$ para os $201$ valores de $n$ verificados, obtidos pela redução e
não por uma raiz.}

\medido{\code{tests/metalica.c} --- $24$ unidades, resíduo $0$, e \textbf{sem um único número de vírgula flutuante no ficheiro}: se um aparecesse, a afirmação de que a família metálica dispensa a reta cairia com ele.}

\noindent \textbf{O índice da família é o traço}, lido dentro do anel; e a norma é $-1$, que é a
Lei~$2$. Não há aqui um número irracional a ser aproximado --- há um anel, uma estaca, e uma
equação entre inteiros.

\begin{proposicao}[o discriminante é o determinante do emparelhamento]\label{prop:tracoform}
A forma $\langle x,y\rangle = \operatorname{tr}(xy)$ é bilinear, e na base $\{1,\sigma\}$ a sua
matriz de Gram é $\begin{psmallmatrix}2&n\\ n&n^{2}+2\end{psmallmatrix}$, cujo determinante é
$\Delta = n^{2}+4$. A forma nunca degenera.
\end{proposicao}

\medido{Verificado para os $101$ primeiros valores de $n$: a Gram é exatamente essa e o
determinante é exatamente $n^{2}+4$, com resíduo zero. E $\Delta\geq 4$ sempre, pelo que a leitura
existe em todas as instâncias da família.}

\noindent Isto merece ser dito devagar, porque arruma uma quantidade que costuma aparecer sem
explicação. \textbf{O discriminante não é uma quantidade acessória da equação: é o determinante do
emparelhamento dual.} Quando se pergunta se $\Delta$ é ou não um quadrado, está a perguntar-se se
a leitura separa os dois lados ou se os colapsa --- e a resposta decide se a instância é um corpo
ou se decompõe.

\medido{$\Delta=n^{2}+4$ é quadrado perfeito \emph{uma única vez} em $100\,001$ valores de $n$
testados por busca inteira: em $n=0$. É o único $n$ em que o anel decompõe.}

\noindent E esse $n=0$ é o degrau zero da escada: aí $\sigma^{2}=1$ e $\sigma^{\dagger}=-\sigma$,
que é a Lei~$1$ escrita no próprio anel. \emph{A involução não é um caso à parte da família
metálica --- é o seu termo inicial}, e para todo $n\geq 1$ a igualdade
$\sigma^{\dagger}=-\sigma$ falha, o que confirma que a Lei~$1$ marca um degrau e não decora todos.

\subsection{A potência fecha, e a torre alterna --- tudo em inteiros}\label{sec:potfecha}

\medido{$\sigma^{k}=F_{k-1}+F_{k}\sigma$, com $F$ a recorrência de índice $n$, verificado por
\emph{dois caminhos independentes} (potência dentro do anel \emph{versus} recorrência) em $132$
casos, sem discordância. O traço fecha e obedece a $t_{k+1}=n\,t_{k}+t_{k-1}$ em $184$ casos. E
$\operatorname{N}(\sigma^{k})=(-1)^{k}$: \textbf{a alternância preto-e-branco da torre é a
multiplicatividade da norma}, e lê-se em inteiros.}

\subsection{E a fronteira: o que é álgebra e o que é medida}\label{sec:fronteira-alg}

\noindent Vale distinguir duas perguntas que se confundem com facilidade, porque a distinção é
exatamente o critério que decide o que fica neste volume e o que desce ao \emph{Catálogo}.

\begin{center}\small
\begin{tabular}{@{}lll@{}}
\toprule
\textbf{a pergunta} & \textbf{do que precisa} & \textbf{onde vive} \\
\midrule
$x$ é unidade? & da norma, que é do anel & \textbf{aqui} --- é álgebra \\
o conjugado é pequeno? & de comparar tamanhos & \textbf{Catálogo} --- é medida \\
\bottomrule
\end{tabular}
\end{center}

\medido{Nas $3\,240$ formas de grau dois varridas, os dois critérios \emph{não} coincidem: há
$1\,522$ que satisfazem o critério métrico sem serem unidades, e casos de unidades que não o
satisfazem. Nenhum dos dois implica o outro. E a família metálica está na interseção, sem exceção
na amostra --- o que é a razão de o teorema valer para ela: não por acaso, por estar dos dois
lados.}

\noindent \emph{É esta a linha de corte de todo o texto.} O que se resolve com a estaca e a cruz
fica; o que exige comparar tamanhos é realização, e vive onde as realizações vivem.


\section{O enunciado}\label{sec:enunciado}

Este texto tem \emph{uma} afirmação. Tudo o resto --- as trinta secções, as vinte e cinco
proposições, os medidores do catálogo --- é realização dela ou consequência sua. Enuncia-se aqui,
inteira, e não volta a enunciar-se.

\begin{mdframed}[linewidth=0.4pt,innerleftmargin=10pt,innerrightmargin=10pt,
                 innertopmargin=8pt,innerbottommargin=8pt,skipabove=6pt,skipbelow=6pt]
\textbf{Um corpo não opera sozinho.} Operar exige \emph{duas} coordenadas ligadas por uma involução
$\nu$ com $\nu\circ\nu=\mathrm{id}$: uma que \textbf{mede} e outra que \textbf{ordena}. A que mede é
aditiva, dá tamanho e não distingue lados; a que ordena é multiplicativa, dá sentido e não dá
tamanho. Nenhuma das duas fecha sozinha, e a operação é o que acontece entre elas.

\medskip
Daí três consequências, e são as únicas de que este texto precisa:
\begin{enumerate}\itemsep2pt
\item \textbf{a fronteira é onde $\nu$ tem ponto fixo} --- o sítio onde as duas coordenadas
coincidem e portanto deixam de separar;
\item \textbf{o que não tem imagem por $\nu$ não pertence} --- não por proibição, por não haver
como voltar;
\item \textbf{subir de dimensão não acrescenta tamanho: acrescenta fase.} O que a dimensão nova
traz é uma \emph{direção de rotação}, não um sítio novo onde pôr grandeza: a coordenada que mede
fica, a que ordena é que muda.
\end{enumerate}
\medskip

\hrule
\medskip
\noindent \textbf{O QUE DAÍ SAI, EM CONCRETO.} \emph{Todo elemento do lado completado é o limite
de uma órbita de operações do lado discreto, e \textbf{cada passo} inverte-se lá dentro.} Não se
reverte um símbolo --- não há notação a decodificar, nem convenção que possa mentir.
\textbf{Reverte-se o objeto}, pela estaca, sem aproximação em nenhum passo:
\[
  x \;\longleftrightarrow\; x^{\dagger},
  \qquad
  M_w = M_{a_0}M_{a_1}\cdots M_{a_{k-1}}, \qquad \det M_w = (-1)^k ,
\]
e o sinal alternado do determinante é a estaca a contar quantas vezes trocou de lado.
\textbf{As duas colunas de $M_w$ são a cruz} --- o par que encaixota o objeto ---, e é por isso
que a reversão é exata e não ortográfica.

\medskip
\noindent \emph{Nomear os conjuntos em que isto acontece é escolher uma instância, e não é
preciso aqui:} a construção completa, com $\N$, $\Z$, $\Q$ e $\R$ levantados um do outro pela
estaca e conservados pela cruz, está no \emph{Catálogo}, que é onde as realizações vivem.

\medskip
\noindent \textbf{E os \emph{pontos fixos} desta operação são exatamente os elementos
quadráticos} --- as órbitas periódicas, por Lagrange. Não é uma restrição de método: se $M\neq\pm
I$ é unimodular e fixa $x$, então $x$ satisfaz uma equação de grau dois com coeficientes no anel
de base, \emph{e mais nada}. O que é transcendente fica de fora \emph{por um teorema}, não por
uma escolha. \emph{Quem alcança tudo é a órbita; o ponto fixo alcança um bocado dela.}

\medskip
\noindent \textbf{E a mesma matriz é a cifra}: $\det=\pm1$ dá ao mesmo tempo a inversa inteira e a
decifra de \emph{todo} criptograma --- com $\det=2$, metade não decifra. \emph{Um facto a servir
dois papéis} (e não uma condição que seja a mesma nos dois: no degrau, $\det=\pm1$ é forçado por
construção; na cifra, pode falhar).

\medskip
\medskip
\noindent \textbf{E o que a Parte III acrescenta, com o número ao lado:}

{\small
\begin{longtable}{@{}p{0.44\textwidth}p{0.52\textwidth}@{}}
\hline
\textbf{a transformada é a avaliação nas raízes} & e para a borda são as \emph{folhas} de Frobenius \\
\quad e a \emph{dourada} é o caso mais simples & Fourier sozinho é \emph{metade}, e não mede este objeto \\
\quad porque as raízes do metal & \textbf{não estão no círculo}: $|\sigma||\sigma'|=1$, recíprocas \\[2pt]
\textbf{a base já é a função} & $28\,561$ pares de $\mathrm{GF}(13^2)$, produto diagonal \\
\quad e o $\Delta$ é o preço da métrica errada & $\det\mathrm{Gram}=m^2+4$, na régua do traço \\[2pt]
\textbf{o completado é um lado somado ao seu dual --- \emph{duas} transformadas} & pares abaixo, ímpares acima, alternância exata \\
\quad cada lado com razão $\sigma^2$ & $x_k = (m^2{+}2)x_{k-1}-x_{k-2}$ \\
\quad e o ouro não salta nó nenhum & corrida $1$ na árvore --- \emph{vazamento zero} \\[2pt]
\textbf{a deconvolução nunca falha na zeta} & porque $\zeta(0)=1/\det(I)=1$ é unidade do anel de séries \\
\quad e falha onde a borda cinde & o núcleo constante anula $15$ de $16$ \\[2pt]
\textbf{$f'=f^{-1}$ força o ouro} & $b^2-b-1=0$ \emph{é} a borda com $m=1$ \\
\hline
\end{longtable}}

\medskip
\noindent \textbf{O que é próprio deste texto, dito sem exagero:} o material está em Khinchin,
Perron, Lagrange e Baire, e a chave matricial unimodular é o cifrador de Hill (1929) --- tudo
clássico, e citado. \emph{Próprio é o enquadramento por $\nu$} --- ler a completação como involução do
par de convergentes, e ver a mesma matriz a servir o degrau e a cifra --- \emph{e a realização
medida e executável}: o degrau no \emph{Catálogo}, \code{tests/palavra.c} ($33$ unidades, resíduo $0$)
e \code{tests/nomeia.c}, que opera sem um único \emph{float}.
\medskip
\end{mdframed}

\paragraph{A lei de que tudo isto sai, e ela é uma exigência sobre a dinâmica.} Antes da
palavra vem o pedido, e o pedido é um só: \textbf{que todo o passo seja reversível}. Não é uma
propriedade que se observe --- é uma condição que se impõe, e o resto é o que sobrevive a ela.

A lei escreve-se numa linha e lê-se por níveis: \emph{o nível é quantas vezes se deriva antes de
inverter.}

\begin{mdframed}[linewidth=0.4pt,innerleftmargin=10pt,innerrightmargin=10pt,
                 innertopmargin=8pt,innerbottommargin=8pt,skipabove=6pt,skipbelow=6pt]
\begin{center}
\begin{tabular}{@{}c@{\quad}l@{\quad}l@{}}
\textbf{nível 0} & $T^{\dagger}=T$ & involução, \emph{dualidade pura}, a Möbius involutiva \\[2pt]
\textbf{nível 1} & $f'=f^{-1}$     & \emph{bidualidade}: a Möbius de Fibonacci, e a razão áurea \\[2pt]
\textbf{nível $n$} & $f^{(n)}=f^{-1}$ & a família metálica inteira \\
\end{tabular}
\end{center}
\end{mdframed}

\medskip
\noindent \textbf{Nível 0 --- a dualidade pura.} $T^{\dagger}=T$ é o pedido mais nu que existe: \emph{ir
é voltar}. Numa transformação de Möbius $f(z)=(az+b)/(cz+d)$ ela tem uma tradução exata e
puramente aritmética:
\[
  f=f^{-1} \iff a+d=0 \quad\text{(traço nulo)} ,
\]
\emph{medido em $6016$ matrizes, zero discordâncias}. É a régua autodual deste texto, e o nível
onde ainda não há metal nenhum: a função é o seu próprio dual e mais nada.

\medskip
\noindent \textbf{Nível 1 --- e é aqui que nasce o ouro.} Pedir agora que \emph{derivar uma vez}
seja inverter obriga, para $f=a\,x^{b}$, a $b-1=1/b$:
\[
  b^{2}-b-1=0 \iff b=\varphi .
\]
E a mesma coisa em Möbius é a matriz de Fibonacci $M=\bigl(\begin{smallmatrix}1&1\\1&0\end{smallmatrix}\bigr)$,
isto é $f(z)=1+1/z$, cujo \emph{ponto fixo} $z=1+1/z$ é $\varphi$. Note-se: $\operatorname{tr}M=1$,
\emph{não} zero --- $M$ não é involução. Ela é o \emph{gerador}.

\medskip
\noindent\emph{Nenhuma das duas equações é nova, e convém dizê-lo já.} A equação $f'=f^{-1}$ está
resolvida na literatura, e a solução é a potência de expoente $\varphi$~\cite{cook2021}; a
fatorização que se segue é um caso de um teorema de 1966--67~\cite{wonenburger1966,djokovic1967}.
\textbf{Elas aparecem separadas, cada uma por si.} O que aqui se faz é pô-las na mesma escada
--- e é a \emph{análise da família de potência sob a lei}, no \emph{Catálogo} (é realização), que mostra o
que nenhuma delas mostra sozinha: que a escada é uma fórmula só, que o coeficiente nunca falha,
que as raízes são o par da alfândega, e que a paridade de~$n$ decide quantas soluções existem.


\medskip
\noindent \textbf{Integrando, a lei dá a conservação.} Para $f$ \emph{crescente} o retângulo
reparte-se em dois e não sobra nada,
\[
  \int_a^{b} f  \;+\; \int_{f(a)}^{f(b)} f^{-1} \;=\; b\,f(b) - a\,f(a) ,
\]
verificado nos três primeiros metais com resíduo $10^{-14}$; para $f$ \emph{decrescente} --- e a
Möbius de Fibonacci é decrescente --- a mesma conta dá \emph{diferença} em vez de soma, com o mesmo
membro direito. \textbf{Isto é Parseval antes de haver transformada}, e reaparece como
$\|x\|^2=\|\mathcal F x\|^2$ (a transformada universal no \emph{Catálogo}).

\medskip
\noindent E no ponto fixo a repartição separa as duas coordenadas deste texto, exatamente:
\[
  \int_1^{\varphi}\Bigl(1+\tfrac1x\Bigr)dx=\underbrace{\tfrac1\varphi}_{\text{algébrica}}+\underbrace{\ln\varphi}_{\text{analítica}},
  \qquad
  \int_{\varphi}^{2}\frac{dy}{y-1}=\ln\varphi ,
  \qquad
  \text{diferença}=\frac1\varphi .
\]
\textbf{A dualidade cancela o transcendente e deixa a borda.} O $\ln\varphi$ é o que as duas
metades \emph{partilham}; $1/\varphi=\varphi-1$ é o que as \emph{distingue} --- e é a borda
$\varphi^{2}=\varphi+1$. Memória da divisão, escrita em integral.

\medskip
\noindent Nas potências $M^{n}=\bigl(\begin{smallmatrix}F_{n+1}&F_n\\F_n&F_{n-1}\end{smallmatrix}\bigr)$ a
divisão dá $M^{n}(x)=F_{n+1}/F_n+(-1)^{n+1}\!/\bigl(F_n(F_nx+F_{n-1})\bigr)$, logo
\[
  \int M^{n} \;=\; \frac{F_{n+1}}{F_n}\,x \;+\; \frac{(-1)^{n+1}}{F_n^{2}}\,\ln\bigl(F_nx+F_{n-1}\bigr) .
\]
\emph{O coeficiente do logaritmo é a identidade de Cassini} $F_{n+1}F_{n-1}-F_n^{2}=(-1)^n$ ---
verificado para $n=2\ldots10$. A parte linear é o convergente $F_{n+1}/F_n$; a parte logarítmica
carrega o sinal de $\det M=-1$. \textbf{A alfândega $\sigma\sigma'=-1$ aparece na integral como o
sinal do logaritmo.}


\medskip
\noindent \emph{Da lei saem, por esta ordem: a involução; o ouro; a família dos corpos; a
conservação; a régua única; e daí a transformada, a convolução, e todas as dualidades que a
literatura já nomeou.} O que vem a seguir é a palavra que descreve o mecanismo.

\paragraph{E antes de tudo, a palavra: o que aqui se chama \emph{dualidade}.} Ela aparece em
todas as partes deste texto e em roupas diferentes --- na álgebra, na geometria, na dimensão e na
lógica --- e por isso fica dita uma vez, inteira:

\begin{mdframed}[linewidth=0.4pt,innerleftmargin=10pt,innerrightmargin=10pt,
                 innertopmargin=8pt,innerbottommargin=8pt,skipabove=6pt,skipbelow=6pt]
\textbf{A dualidade é a memória da divisão.}

\medskip
Dividir é uma operação que \emph{perde}. Parta-se um objeto em duas partes e guarde-se uma: das
duas metades ficou uma, e da que ficou não se reconstrói o que era --- medido, num universo de
$169$ elementos, guardar só a parte simétrica colapsa-o em $37$ e \textbf{perdem-se $132$
distinções}; guardar só a antissimétrica perde $156$. \emph{Guardando as duas, perde-se zero.}

\medskip
\textbf{A dualidade é o que guarda a segunda} --- a memória de que as duas metades eram uma. É ela
que torna a divisão \emph{reversível}, e por isso é uma tradução biunívoca: uma bijeção $\nu$ com
$\nu\circ\nu=\mathrm{id}$, em que \emph{a volta é a própria ida} e não há um segundo mapa a
inventar para regressar.

\medskip
\textbf{A sua realização é a antissimetria.} Toda a dualidade parte o objeto em duas:
\[
  x \;=\; S + A, \qquad \nu(S)=S, \qquad \nu(A)=-A .
\]
A parte \emph{simétrica} é o que se guarda; a \emph{antissimétrica} é o que executa a troca.
\emph{Sozinha, a simétrica não reverte nada} --- ela não distingue $x$ de $\nu(x)$. Reverter é
trocar o sinal de $A$ e manter $S$: um passo, e não um caminho de regresso.

\medskip
\textbf{E ela fecha porque é uma graduação por $\{\pm1\}$:}
\[
  S\cdot S \subseteq S, \qquad S\cdot A \subseteq A, \qquad \boxed{A\cdot A \subseteq S\,}
\]
A última é a que diz tudo: \emph{o produto de dois antissimétricos é simétrico} --- a involução
torna a antissimetria simétrica, e é por isso que a dualidade fecha em vez de fugir para fora.

\medskip
\textbf{E daqui sai o critério, que é o que separa um par dual de um par qualquer:} se não houver
involução, não há memória, e sem memória não há dualidade --- há apenas duas coisas ao lado uma da
outra. \emph{Este texto já o aplicava antes de o enunciar:} recusa o par $(D,\int)$ como dualidade
porque $\int\!\circ D = \mathrm{id}-\ker D$ e o núcleo é a parte que se perdeu; e recusa a
transformada de Legendre pela mesma razão --- ela é um limite, não um isomorfismo, e não tem volta.
\emph{Onde a divisão perde, não é dualidade: é degeneração.}
\end{mdframed}

\noindent \textbf{As sete roupas, e a mesma pessoa dentro.} Este texto usa as sete; ficam aqui
lado a lado para que se veja que são uma (\code{tests/pontofixo.c}, $40$ unidades, resíduo $0$):

\smallskip
\begin{center}\small
\begin{tabular}{@{}lp{0.30\textwidth}p{0.34\textwidth}@{}}
\hline
& \textbf{a involução} & \textbf{troca / guarda} \\
\hline
\textbf{algébrica} & $\nu(a+b\sigma)$, com $\nu\circ\nu=\mathrm{id}$ & troca as duas coordenadas; guarda a \emph{norma} \\
\textbf{geométrica} & o \emph{cone dual}: $K^{**}=K$ & troca dentro e fora; guarda a \emph{forma} \\
\textbf{dimensional} & $V-E+F=2$, e o dual troca $V$ e $F$ & troca vértices e faces; guarda a \emph{aresta} \\
\textbf{lógica} & a negação: $\neg\neg p = p$ & troca $\wedge$ e $\vee$ (De Morgan); guarda o \emph{valor} \\
\textbf{projetiva} & ponto $\leftrightarrow$ reta em $\mathrm{PG}(2,q)$ & troca ponto e reta; guarda a \emph{incidência} \\
\textbf{de Gelfand} & álgebra $\leftrightarrow$ espectro (as raízes) & troca álgebra e pontos; guarda o \emph{produto} \\
\textbf{bidualidade} & a dualidade sobre si: $V^{**}=V$ & \emph{dois níveis}, e por isso quatro componentes \\
\hline
\end{tabular}
\end{center}


\smallskip
\noindent \textbf{E a projetiva é a mais antiga das seis, e a que lhes dá o nome.} No plano
projetivo, trocar \emph{ponto} por \emph{reta} leva teorema verdadeiro em teorema verdadeiro ---
``dois pontos determinam uma reta'' torna-se ``duas retas determinam um ponto'', Pascal torna-se
Brianchon, e Desargues é \emph{autodual}, porque o seu dual é a própria recíproca. O que sustenta o
princípio é aritmético e conta-se: em $\mathrm{PG}(2,q)$ há $q^2+q+1$ pontos e \emph{exatamente} o
mesmo número de retas, com $q+1$ de cada em cada, e a incidência é simétrica. \emph{Não é analogia
com as outras cinco: é o mesmo mecanismo --- uma bijeção que troca dois papéis e deixa um
invariante de pé.}

\smallskip
\noindent \textbf{E a de Gelfand é a que este texto usa na análise sem lhe chamar o nome.} A
transformada deste projeto \emph{é} a avaliação nas raízes da borda --- e isso é exatamente a
transformada de Gelfand: o espectro de $\Z_p[x]/(f)$ são as raízes de $f$, avaliar em cada uma é o
mapa $a\mapsto(a(\sigma_1),\dots,a(\sigma_n))$, e o teorema diz que ele é isomorfismo quando $f$
cinde. Medido: bijeção sobre os $121$ elementos de $\Z_{11}[x]/(x^2-x-1)$ e
$\widehat{ab}=\hat a\,\hat b$ nos $14\,641$ pares, resíduo $0$. \emph{A álgebra lê-se no espectro e
o espectro reconstrói a álgebra} --- e é por isso que ``a transformada é a avaliação nas raízes''
não é uma escolha de método: é a dualidade de Gelfand no caso finito.

\paragraph{E daí saem a transformada e a convolução, sem se escolher nenhuma.} Se o passo tem de
ser reversível e a régua tem de ser a mesma nos dois sentidos, então a única coisa que se pode
pedir a uma \emph{transformada} é que \textbf{leve a operação composta em operação simples e tenha
volta}. Não se escolhe a fórmula: pede-se a propriedade, e a fórmula vem.

\smallskip
\noindent \textbf{A transformada é a avaliação nas raízes da borda.} Multiplicar em $\Z_p[x]/(f)$ é
convoluir os coeficientes com redução por $f$; avaliar num zero de $f$ é um homomorfismo; logo
avaliar nos zeros leva o produto no produto:
\[
  \mathcal F(a\circledast b) \;=\; \mathcal F(a)\,\mathcal F(b) .
\]
Medido no anel \emph{inteiro}, $14\,641$ pares, resíduo $0$ (a transformada universal no \emph{Catálogo}). E isto é a
dualidade de Gelfand no caso finito --- a álgebra lê-se no espectro, o espectro reconstrói a
álgebra.

\smallskip
\noindent \textbf{E a convolução é a operação que ela simplifica; a deconvolução, a volta.} Fundir
dois corpos é convoluir; do outro lado é multiplicar casa a casa; e por isso \emph{desfazer é
dividir}. A volta existe exatamente onde a lei o permite: \textbf{sse nenhuma folha se anula} --- e
isso não é só decidível, é contável: dos $121$ núcleos de $\Z_{11}[x]/(f)$, exatamente
$100=(p-1)^n$ têm inversa. \emph{O que anula uma folha é o que não tem dual}, e é o único sítio
onde a reversibilidade falha --- por isso é o único que a lei exclui.

\paragraph{E onde cada uma já vive, para quem for procurá-las.} As dualidades clássicas não
entram aqui como ornamento: o texto já as usava, e este mapa diz onde --- \emph{são os pontos fixos
deste livro}, e quem os conhecer da literatura reconhece a casa.

\smallskip
\begin{center}\small
\begin{tabular}{@{}lp{0.30\textwidth}p{0.32\textwidth}@{}}
\hline
\textbf{a dualidade} & \textbf{onde vive} & \textbf{o que a divisão guarda} \\
\hline
Pontryagin & o eixo de Pontryagin no \emph{Catálogo} --- e é \emph{o eixo} das três partes & $G\cong\widehat{\widehat G}$; compacto\,$\leftrightarrow$\,discreto, com o autodual no ponto fixo \\
Poincaré & Euler e o miolo no \emph{Catálogo}, a escada $\chi=2\to1\to0$ & o emparelhamento entre os graus $k$ e $n-k$ \\
Gelfand & no \emph{Catálogo} --- \emph{a transformada é a avaliação nas raízes} & o produto: $\widehat{ab}=\hat a\hat b$ \\
convolução & no \emph{Catálogo}, e a deconvolução ao lado & $\ast\leftrightarrow\cdot$; e a volta existe \emph{sse} nenhuma folha se anula \\
Fourier\,/\,Mellin & a transformada dourada no \emph{Catálogo} --- ``Mellin é Fourier no multiplicativo'' & o caractere; o $\log$ é o isomorfismo \\
série\,/\,traços & os quatro na zeta dinâmica no \emph{Catálogo} --- Fibonacci contra Lucas & $\zeta(0)=\det(I)^{-1}=1$, e por isso a deconvolução nunca falha \\
projetiva & a tabela acima & a incidência: $q^2+q+1$ de cada em $\mathrm{PG}(2,q)$ \\
\hline
\end{tabular}
\end{center}

\smallskip
\noindent \textbf{E há duas que o texto usa sem lhes chamar o nome, e fica dito.} A representação
de \emph{Riesz} --- identificar o funcional com o vetor --- é o que a forma-traço $\langle
x,y\rangle=\operatorname{tr}(xy)$ faz em toda a Parte~I, e o nome nunca aparece. E o par
\emph{polos\,$\leftrightarrow$\,zeros}: o texto tem os dois lados (os polos de $\zeta$ são os
inversos das raízes da borda) sem os escrever como par involutivo. \emph{Dizer onde a casa está
vazia é parte de mostrar a casa.}

\paragraph{E a terceira consequência é a que arruma as outras duas.} Multiplicar por $i$
preserva o módulo \emph{exatamente} e roda $90^\circ$ --- verificado nas quatro potências: $1$,
$i$, $-1$, $-i$, todas de módulo $1$, com a fase a andar $0^\circ \to 90^\circ \to 180^\circ \to
-90^\circ$. \textbf{O tamanho nunca muda; a ordem é que muda.} E é a mesma repartição que aparece
em cada realização deste texto: a conservação fixa o módulo e deixa a fase livre; onde a fase é $0$
ou $\pi$ há espelho, onde é $\pm\pi/2$ há rotação (Obs.~o $\Gamma$ real no \emph{Catálogo}).

\noindent \emph{Donde ``dimensão'' não é uma quantidade de espaço --- é uma quantidade de fases
disponíveis}, e a projeção para a dimensão de baixo é o esquecimento de uma delas.

\paragraph{O que conta como realização, e o que não conta.} Uma realização exibe as duas
coordenadas, a involução que as troca, e o ponto fixo dela --- e mede o resíduo de
$\nu\circ\nu=\mathrm{id}$. \emph{Exibir uma coordenada e chamar-lhe dual não é realização}: é meia
estrutura com o nome do par (Obs.~a convergência dual no \emph{Catálogo}).

\subsection{A bidualidade de cada face, e a condição sob que ela fecha}\label{sec:bidualfaces}

As sete faces estavam medidas como involuções. Falta o segundo andar --- e o que ele mostra não é a
bidualidade em si, que é sempre a mesma involução, mas \textbf{o que a fecha em cada uma}. Note-se
bem: \emph{a bidualidade fecha sempre}. A coluna da direita não lista condições que a limitem ---
lista \textbf{a estrutura que a realiza} em cada face.

\begin{center}\small
\begin{tabular}{@{}llll@{}}
\toprule
face & o dual & o bidual & \textbf{o que a fecha} \\
\midrule
lógica       & $\neg p$              & $\neg\neg p=p$          & o valor de verdade \\
Poincaré     & $H_{n-k}$             & $H^{k}$                 & a classe fundamental \\
dimensional  & o poliedro dual       & o próprio               & a aresta, que $\chi=2$ conserva \\
projetiva    & a recta               & o ponto                 & a incidência \\
geométrica   & o cone $K^{*}$        & $K^{**}=K$              & a convexidade (o bipolar) \\
Pontryagin   & $\widehat G$          & $\widehat{\widehat G}\cong G$ & a medida de Haar \\
Gelfand      & o espectro            & a álgebra               & o produto \\
algébrica    & $A^{*}$               & $A\times A^{*}$         & a norma --- e a dimensão dobra \\
\bottomrule
\end{tabular}
\end{center}

\medido{Lógica nos $2$ valores e De Morgan nos $4$ pares; Poincaré em $860$ pares $(n,k)$;
dimensional nos $5$ sólidos platónicos, com $\chi=2$ e o bidual a devolver o próprio; projetiva em
$q=2,3,5,7$ com $\#\text{pontos}=\#\text{rectas}=q^{2}+q+1$; Pontryagin nos cíclicos de ordem $n$, $n=2\ldots24$,
onde $|\widehat G|=|G|$ e o bidual devolve $G$. Tudo em \code{tests/bidual.c}, resíduo $0$.}

\begin{obs}[e o que fecha é sempre o mesmo, noutra roupa]
O valor, a classe fundamental, a aresta, a incidência, a convexidade, Haar, o produto, a norma:
\textbf{todos são \emph{o invariante} --- a coisa que a troca não move.} Dualizar é trocar dois
papéis e deixar um terceiro quieto, e é o terceiro que faz a volta ser possível. Por isso a
bidualidade não precisa de licença para valer: \emph{ela fecha porque há sempre alguma coisa que não
se moveu}.

\smallskip
\noindent O que \emph{varia} é quanto do objeto o invariante alcança. Num espaço vetorial de dimensão
finita ele alcança tudo e $V^{**}\cong V$; em dimensão infinita a injeção canónica $V\to V^{**}$ não é
sobrejetiva e o bidual é \emph{estritamente maior} --- a volta dá-se na mesma, mas chega a um sítio
maior do que aquele de onde se partiu.

\smallskip
\noindent \textbf{É por isso que o corpo deste texto é finito}: não porque a dualidade precise disso
para existir, mas porque é aí que o invariante cobre o objeto inteiro e ler e escrever usam
\emph{exatamente} a mesma régua.
\end{obs}

\subsection{A existência do dual, e as duas leis}\label{sec:duasleis}

Há dois andares, e convém não os confundir. O primeiro é a \textbf{existência}: que o dual há de
existir, e sem que se lhe peça nada. O segundo são \textbf{as duas leis}, que dizem \emph{que forma}
ele tem. Ambos se provam a partir das mesmas duas peças que este texto usa desde o princípio: a
involução e o dicionário.

\paragraph{Primeiro andar: a existência.} Três proposições, e nenhuma usa hipótese sobre o objeto.

\begin{proposicao}[a escrita tem dual]\label{thm:lei1}
Seja $\rho\colon G\to\GL(V)$ uma representação. Então
\[
  \rho^{*}(g) \;:=\; \bigl(\rho(g)^{-1}\bigr)^{\!\top}
\]
é também uma representação de $G$, e $(\rho^{*})^{*}=\rho$. \emph{A dualidade não é uma propriedade
que algumas estruturas tenham: é uma construção que existe sempre.}
\end{proposicao}

\begin{proof}
Homomorfismo: $\rho^{*}(gh)=\bigl((\rho(g)\rho(h))^{-1}\bigr)^{\top}
=\bigl(\rho(h)^{-1}\rho(g)^{-1}\bigr)^{\top}
=\bigl(\rho(g)^{-1}\bigr)^{\top}\bigl(\rho(h)^{-1}\bigr)^{\top}=\rho^{*}(g)\,\rho^{*}(h)$ --- a
transposição inverte a ordem, e a inversão inverte-a outra vez; as duas trocas cancelam-se.
Involução: como $(A^{\top})^{-1}=(A^{-1})^{\top}$, tem-se
$\bigl((M^{-1})^{\top}\bigr)^{-1}=M^{\top}$, donde $\nu(\nu(M))=(M^{\top})^{\top}=M$.
Nenhum dos dois passos usa hipótese alguma sobre $G$ ou $V$.
\end{proof}

\medido{$\rho^{*}(gh)=\rho^{*}(g)\rho^{*}(h)$ em $14\,400$ pares de $\GL_2$ e
$\nu\circ\nu=\mathrm{id}$ nas $232$ matrizes de $|\det|=1$ com entradas em $[-3,3]$: \textbf{zero
falhas} em ambos (\code{tests/bidual.c}).}

\paragraph{E é aqui que entra o dicionário.} A construção existe sempre, mas para ela ficar
\emph{dentro do anel} é preciso que $\nu(M)=(M^{-1})^{\top}=\bigl(\operatorname{adj}M\bigr)^{\top}\!/\det M$
não tenha denominador --- isto é, $|\det M|=1$. E $|\det|=1$ é precisamente o que o dicionário
palavra\,$\to$\,matriz já dá por construção (o dicionário palavra$\to$matriz no \emph{Catálogo}): \emph{a condição do dual no anel de base e a
propriedade do dicionário são a mesma}. Fora dela o dual continua a existir, mas sai do anel de base: para
$\bigl(\begin{smallmatrix}2&0\\0&1\end{smallmatrix}\bigr)$, de determinante $2$, ele tem entrada
$1/2$.

\begin{proposicao}[a leitura tem dual]\label{thm:lei2}
Seja $\mu$ uma \emph{leitura} --- uma medida que a cada objeto associa um escalar. Então
$\mu^{*}(x):=\mu(\nu(x))$ é também uma leitura, e $\mu^{**}=\mu$.
\end{proposicao}

\begin{proof}
$\mu^{*}$ leva a mesma classe nos escalares, logo é leitura; e
$\mu^{**}(x)=\mu^{*}(\nu(x))=\mu(\nu(\nu(x)))=\mu(x)$ pela Primeira Lei.
\end{proof}

\noindent \emph{A leitura matemática desta é a integral} --- e é ela que \textbf{impede a perda}.
Guardar uma leitura sozinha colapsa o objeto; guardar o par não perde nada:

\begin{center}\small
\begin{tabular}{@{}lll@{}}
\toprule
a leitura & a sua dual & e o que as liga \\
\midrule
traço $\sigma+\sigma'$ & determinante $\sigma\sigma'$ & soma $\leftrightarrow$ produto (Vieta) \\
dimensão $k$           & codimensão $n-k$             & $k+(n-k)=n$ (Poincaré) \\
vértices $V$           & faces $F$                    & $V-E+F=2$ (Euler) \\
norma $|\sigma|$       & norma $|\sigma'|$            & $|\sigma||\sigma'|=1$ (a alfândega) \\
$D f$                  & $f$ na fronteira             & $\int\!Df = f - f(a)$ \\
\bottomrule
\end{tabular}
\end{center}

\medido{$124$ bordas distintas colapsam em $13$ traços --- $111$ distinções perdidas com a leitura
sozinha, e \emph{zero} com o par. E na última linha: $\int\!D$ sem guardar a fronteira falha em
$392$ de $441$ polinómios; guardando o par $(Df,\,f(0))$, falha em \textbf{zero}. \emph{É por isto
que $(D,\int)$ não é dualidade sozinho} --- e é o que o texto já recusava, agora com a razão dita:
falta-lhe o segundo membro, que é a fronteira (\code{tests/escada.c}).}

\begin{proposicao}[leitura e escrita são adjuntas]\label{thm:lei3}
Escrever com $A$ e depois ler com $y$ é o mesmo que ler com $A^{\top}y$ e depois escrever:
\[
  \langle A x,\, y\rangle \;=\; \langle x,\, A^{\top} y\rangle .
\]
\emph{Mover a ação de um lado para o outro troca $A$ por $A^{\top}$.}
\end{proposicao}

\medido{$1200$ casos em dimensões $2$, $3$ e $4$, com entradas inteiras: resíduo $0$.}

\noindent \textbf{E esta contém as duas primeiras.} Fixando $y$, o lado esquerdo é uma leitura e a
identidade exibe a sua dual --- é a segunda proposição. Fixando $x$, obtém-se a primeira por
$A\mapsto A^{\top}$. \emph{As três não são independentes: a terceira é o par, e as outras duas são
ela vista de cada lado.}

\paragraph{E é a terceira lei de Newton.} \emph{A toda ação corresponde uma reação igual e
oposta}; aqui, a toda escrita corresponde uma leitura igual e oposta. A correspondência é literal e
não analógica: em Möbius, a adjunção na forma $T^{\dagger}=T$ equivale a traço nulo, isto é
$d=-a$ --- \textbf{a segunda entrada da diagonal é a primeira com o sinal trocado}, exatamente como
$F_{AB}=-F_{BA}$. \emph{Medido: a equivalência ``traço nulo $\iff d=-a$'' em $13\,808$ matrizes,
zero discordâncias.}

\paragraph{Segundo andar: as duas leis, e as provas.} A existência não diz nada sobre a
\emph{forma} do dual. É isso que as duas leis fixam, e cada uma prova-se numa linha.

\begin{teorema}[Lei 1 --- a unidade é dual]\label{thm:leiA}
As unidades do anel formam um par trocado pela involução, e não um elemento sozinho:
$1^{\dagger}=-1$, isto é, têm \emph{o mesmo tamanho e sinais opostos}. Numa transformação de Möbius isso
equivale a
\[
  \operatorname{tr} = 0, \quad\text{isto é}\quad d=-a .
\]
\end{teorema}

\begin{proof}
As unidades de $\Z$ são as soluções de $x^{2}=1$, ou seja $\{+1,-1\}$: dois elementos de módulo $1$
e sinais opostos, e a involução troca-os. Em $\GL_2$, $M^{2}=\lambda I$ com $M\neq\pm I$ obriga
$b(a+d)=c(a+d)=0$ com $b,c$ não ambos nulos, logo $a+d=0$; e reciprocamente traço nulo dá
$M^{2}=(a^{2}+bc)I$, escalar --- que em Möbius \emph{é} a identidade.
\end{proof}

\medido{A equivalência ``traço nulo $\iff d=-a$'' em $13\,808$ matrizes invertíveis, zero
discordâncias (\code{tests/bidual.c}).}

\begin{teorema}[Lei 2 --- a dualidade é dual]\label{thm:leiB}
A própria involução tem dual, e a forma disso é $T^{\dagger}=-T$: \emph{o dual de $f$ é o simétrico da
sua inversa}. Sobre a raiz da borda lê-se exatamente
\[
  \sigma' = -\sigma^{-1} \iff \sigma\sigma' = -1 \iff \det = -1 ,
\]
que é a assinatura da Möbius de Fibonacci. E aplicá-la duas vezes devolve o próprio.
\end{teorema}

\begin{proof}
Por Vieta, o produto das raízes de $x^{2}-mx-1$ é o termo independente, $-1$; logo
$\sigma'=-1/\sigma$. Reciprocamente, $\sigma\sigma'=-1$ força o termo independente a $-1$, isto é
$\det=-1$. E a bidualidade: aplicar $\nu$ duas vezes dá $\nu(\nu(\sigma))=\nu(\sigma')=\sigma$,
porque $\nu$ troca as duas raízes e mais nada.
\end{proof}

\medido{$\sigma'=-1/\sigma$ em $m=-9\ldots9$ com resíduo $0$; e $\nu\circ\nu=\mathrm{id}$ nas $232$
matrizes de $|\det|=1$ com entradas em $[-3,3]$ (\code{tests/bidual.c}).}

\paragraph{E as duas são os dois objetos que este texto usa desde o princípio.} A Lei~1 dá a Möbius
\textbf{involutiva} --- traço $0$, período $2$, \emph{o espelho que mede}; a Lei~2 dá a de
\textbf{Fibonacci} --- $\det=-1$, \emph{o gerador que anda}. \emph{E é o produto delas que dá o
passo}: $M=J\cdot i$, medido acima.

\paragraph{E daí a família metálica, como consequência e não como lei.} O grau $n$ em que
$f^{(n)}=f^{-1}$ é uma leitura do passo, logo tem dual. Para $f=a\,x^{b}$ existe para todo $b\neq0$
--- \emph{não é exigido ao passo, é medido nele}:
\[
  n=b-\tfrac1b=\sigma+\sigma'=\operatorname{tr},
  \qquad\text{e a sua dual}\qquad
  N=\sigma\sigma'=-1 .
\]
Ter as duas \emph{é} ter a borda $x^{2}-nx-1$. E $n=0$ é o passo que já coincide com o seu dual: o
mesmo traço nulo da adjunção, agora por uma razão e não por coincidência.

\medido{$n=\sigma+\sigma'=\sigma-1/\sigma$ em $17$ casos, resíduo $0$ (\code{tests/escada.c}).}

\noindent \textbf{E as duas leituras não colidem: são o eixo de Pontryagin.} A mesma exigência dá um
período no lado discreto e uma borda irracional no lado contínuo, e em grau dois elas nunca se
encontram --- $\sigma_n>1$ é real, logo nunca está sobre a circunferência, logo nunca tem ordem
finita. É preciso subir de grau para que um mesmo objeto tenha as duas leituras
(onde Pisot falha no \emph{Catálogo}).

\noindent O ciclo fecha: parte-se de $\N$, sobe-se por involuções até $\R^n$, colhem-se os traços
(que são \emph{inteiros}), soma-se em p.u., e a zeta devolve os $\sigma$ de onde tudo saiu.


\paragraph{Como ler o resto.} Cada secção traz o seu estatuto ao lado do título, e são três:

\begin{center}\small
\begin{tabular}{ll}
\hline
\emph{[consequência do enunciado]} & 11 secções --- sai da caixa acima, e não acrescenta hipótese \\
\emph{[realização]} & 15 --- exibe as duas coordenadas, a involução e o resíduo \\
\emph{[ferramenta]} & 3 --- notação e maquinaria, sem afirmação própria \\
\emph{[navegação]} & 3 --- situa o leitor, não afirma \\
\hline
\end{tabular}
\end{center}

\noindent \emph{Uma secção que não seja nenhuma das três não devia estar aqui} --- e o critério é
verificável: procure-se, em cada realização, o par nomeado e o número medido. Onde não estiverem, a
secção está a \emph{ilustrar} em vez de realizar.

\noindent \textbf{Corrido sobre as $15$ realizações: passam as $15$.} E o teste \emph{pode} falhar
--- aplicado às secções que \emph{não} são realizações, falha em seis: o dicionário e os
convergentes (ferramentas, sem par nem medida), a introdução e a conclusão, e ``De volta à
projeção'' (consequência, sem número). \emph{Sem esse controlo negativo, ``15 de 15'' não valeria
nada}: seria uma condição que tudo satisfaz.

\paragraph{E o que este texto \emph{não} afirma.} Que a afirmação acima seja necessária, ou única,
ou que valha fora do que se mediu. Ela é uma \textbf{escolha de estrutura} --- paga-se por render, e
o que aqui se documenta é o que rendeu. Onde uma consequência foi \emph{obrigada} por uma hipótese
independente, isso vai dito no sítio (a involução gera a dimensão no \emph{Catálogo}); onde não foi, vai dito também.

\part{Topologia --- a torre, e a ordem que sobe}
\noindent\hfill\emph{\small[\textbf{realização} da Lei~1: a ordem, e como ela passa de andar em andar]}\medskip

\subsubsection*{E a ordem sobe pela torre --- é a involução que a faz subir}

A ordenação de $\Q(\sigma)$ vale em dimensão dois. \emph{E acima?} Sobe por indução, e não é preciso
inventar nada: \textbf{cada andar é dual do anterior}, e a involução alterna o sinal.

\begin{center}\small
\begin{tabular}{@{}crrrl@{}}
\toprule
andar & corpo & $\dim$ & $(-1)^{n}$ & \\
\midrule
$0$ & $\R$ & $1$  & $+$ & branco \\
$1$ & $\C$ & $2$  & $-$ & \textbf{preto} \\
$2$ & $\mathbb{H}$ & $4$  & $+$ & branco \\
$3$ & $\mathbb{O}$ & $8$  & $-$ & \textbf{preto} \\
\bottomrule
\end{tabular}
\end{center}

\noindent \textbf{A torre é preta e branca}, e o sinal que a alterna é o mesmo $(-1)^{n}$ de
Cassini, do $\det(M^{n})$ e da dicotomia par/ímpar. \emph{É o mesmo corpo em todos os andares} ---
só a dimensão muda, e ela conta quantas vezes se dualizou (no \emph{Catálogo}).

\paragraph{E a ordem continua a do andar de baixo.} Em $A\times A^{*}$ ordena-se primeiro pela
componente que vem de $A$, e só em caso de empate pela de $A^{*}$: \emph{se a ordem de $A$ é total,
a de $A\times A^{*}$ também é}. A indução fecha, e é a involução que a conduz --- \textbf{ela faz
isso naturalmente}, sem convenção acrescentada.

\medido{Totalidade e antissimetria em $625$ pares, zero falhas; e compatibilidade com a soma em
$2401$ comparações, também zero (\code{tests/bidual.c}).}

\paragraph{E onde fica o zero? Há dois, e o logaritmo liga-os.} A torre começa em $\R$, dimensão
$1$. \emph{Abaixo dela está o espaço trivial} $\{0\}$, de dimensão $0$ --- e ele é o \textbf{único
que é o seu próprio dual sem precisar de ponte}: $\operatorname{Hom}(\{0\},K)=\{0\}$. A torre não o
constrói: \emph{parte dele}, e ele é o zero de cada andar.

\noindent E há um segundo zero, do outro lado do par. No aditivo o neutro é $0$; no
multiplicativo é $1$; e $\log 1=0$. Portanto \textbf{$|r|=1$ é o zero do crescimento} --- não cresce
nem encolhe ---, e é exatamente o que acontece em $n\equiv5\ (\mathrm{mod}\ 6)$, onde as raízes do
sexto ciclotómico ficam \emph{sobre} a circunferência. \emph{Aí, dimensão zero quer dizer zero de
expansão.}

\paragraph{E nesse ponto os símbolos colidem --- perde-se a sobrejetividade.} Fora do ponto fixo o
par separa-se \emph{por tamanho}, $|\sigma|\neq|\sigma'|$. No ponto fixo os dois têm módulo $1$ e o
tamanho já não distingue: \textbf{a norma colapsa}.

\medido{Em $\GF(p^{2})$, os elementos de norma $1$ contra os não nulos: $6$ de $24$, $8$ de $48$,
$12$ de $120$, $14$ de $168$ --- colisões de $4{:}1$ a $12{:}1$ (\code{tests/bidual.c}).}

\noindent E num conjunto \emph{finito} isso não fica pela injetividade: \textbf{injetiva $\iff$
sobrejetiva}, logo duas entradas gastas numa saída deixam \emph{outra saída por atingir}.

\medido{A equivalência verificada em $44$ aplicações $x\mapsto kx$ em $\Z/n$, sem falha.}

\begin{obs}[e por isso a situação é \emph{origem}, e não \emph{passagem}]\label{obs:origem}
Uma \textbf{passagem} atravessa-se: há antes e há depois. Uma \textbf{origem} só se deixa: não há de
onde vir. No ponto fixo $\nu(r)=r$, logo a involução \emph{não leva a lado nenhum}; e como a norma
colide, também \emph{não há preimagem única} para voltar atrás. \textbf{O ponto não está no meio do
caminho: é onde o caminho começa.}

\smallskip
\noindent \emph{E isto dá mecanismo a uma frase que já estava no texto} --- ``esse ponto é o começo,
então como pode ser o fim?'' (\S\ref{sec:enunciado}). A resposta era, até aqui, uma boa pergunta;
agora tem razão: \emph{ele é origem porque lá os símbolos colidem, e onde há colisão não há
travessia reversível --- só partida.}
\end{obs}

\begin{obs}[e o alcance, que convém dizer]\label{obs:ordemtorre}
A ordem da torre é \textbf{total} e \textbf{compatível com a soma} --- é ordem de \emph{grupo}.
\emph{Não} é compatível com o produto, logo \textbf{não} faz de $\C$ um corpo ordenado: medido,
$144$ falhas em $576$ pares positivos, e o contraexemplo é o esperado --- $i\cdot i=-1$. \emph{E tem
de ser assim}: um corpo com raiz de $-1$ não se ordena, e isso não é remediável por construção
nenhuma. \textbf{A torre ordena o suficiente para ler e comparar} --- é a Lei~$1$ a funcionar em
cada andar --- \emph{e não promete mais do que isso}.
\end{obs}

\medskip
\noindent \textbf{É por isto que a bidualidade fecha sem precisar de sair}: o dual do corpo não é um
corpo novo --- é o mesmo, lido do outro lado. \emph{A estrutura dual deixa de ser convenção
notacional e passa a ser o ambiente em que a lei tem significado.}

\medskip
\noindent \textbf{As duas leis estão em lados opostos, e é isso que as torna não triviais.} A
Primeira diz que o objeto mais simples que existe já é um par: $1$ vive em $V$, $1^{*}$ vive em
$V^{*}$, e a involução troca-os. \emph{Não é que $1$ seja igual a $-1$} --- é que a unidade não vive
sozinha, e as duas metades são \textbf{uma a projeção da outra}. Em Möbius isso é traço nulo,
$d=-a$.

\medskip
\noindent A Segunda diz que a própria dualidade tem dual. Em símbolos, $T^{*}=-T$: \emph{o adjunto
de $T$ é o simétrico de $T$}, sob a identificação acima. Sobre a raiz lê-se exatamente
\[
  \sigma' \;=\; -\frac1\sigma
  \qquad\text{isto é}\qquad
  \boxed{\ \sigma\,\sigma' = -1\ }
\]
--- que é $\det=-1$, \textbf{a assinatura da Möbius de Fibonacci}. E aplicá-la duas vezes devolve o
próprio: a dualidade da dualidade fecha, e daí a órbita de período $4$ que não dissipa.

\medskip
\noindent \emph{As duas são os dois objetos que este texto usa desde o princípio:} a involutiva, que
\textbf{mede} --- o espelho, período $2$ ---, e a de Fibonacci, que \textbf{anda} --- o gerador. E é
o produto delas que dá o passo (\S\ref{sec:duasleis}).

\medskip
\noindent \textbf{E tudo o resto é corolário.} Delas saem, por esta ordem, as leis do movimento e
as quatro equações elementares --- que são os quatro casos de sinal:

\begin{center}\small
\setlength\tabcolsep{4pt}
\begin{tabular}{@{}llp{0.46\textwidth}@{}}
\toprule
o corolário & o que é & de onde vem \\
\midrule
Newton~1 & inércia                & o tamanho não muda sozinho --- Lei~1 \\
Newton~2 & $\int\!F\,dt=\Delta p$ & integrar a Lei~2 dá o que se conserva \\
Newton~3 & $F_{AB}=-F_{BA}$       & o sinal oposto da Lei~1 \\
Newton~4 & gravitação, $1/r^{2}$  & a Lei~2 na forma nula, $\nabla^{2}\varphi=0$ \\
\midrule
$T=1$              & ponto fixo de valor & Lei~1 com sinal $+$ \\
$T=-1$             & ponto fixo de valor & Lei~1 com sinal $-$ \\
$T^{\dagger}=T$    & involução           & Lei~2 com sinal $+$: $T^{2}=+1$ \\
$T^{\dagger}=-T$   & antissimetria       & Lei~2 com sinal $-$: $T^{2}=-1$ \\
\bottomrule
\end{tabular}
\end{center}

\noindent \emph{E a existência é anterior às duas, e está provada} (\S\ref{sec:duasleis}): toda
representação tem dual e toda medida tem dual, sem hipótese nenhuma sobre o objeto. As duas Leis não
afirmam que o dual existe --- \textbf{dizem que forma ele tem}.

\medido{$\sigma\sigma'=-1$ e $|\sigma||\sigma'|=1$ em $m=-9\ldots9$ sem falha; e $f^{2}=-1$ dá
ordem $4$ exata (\code{tests/bidual.c}).}

\medskip
\noindent \textbf{E aqui fica delimitado o estatuto de tudo o que se segue, porque ele mudou.} As
transformações de Möbius foram o ponto de partida \emph{histórico} deste trabalho, e por isso
aparecem em toda a parte. \emph{Deixaram de ser a fonte.} A fonte são o corpo dual e as duas leis
acima; Möbius é \textbf{derivado} --- e o que ele continua a dar, e que nenhuma outra realização dá
tão bem, é a \textbf{representação em fracções contínuas}.

\begin{center}\small
\begin{tabular}{@{}lll@{}}
\toprule
 & \textbf{o que é} & \textbf{estatuto} \\
\midrule
o corpo dual $(V,V^{*},J)$ e as duas leis & a \emph{fonte} & \textbf{fundamento} \\
Álgebra, Topologia, Análise & realizações & \textbf{consequência} \\
a Möbius de Fibonacci, o gato, os convergentes & a \emph{representação} & \textbf{ferramenta} \\
\bottomrule
\end{tabular}
\end{center}

\noindent \emph{E a transição tem um lugar exato}: até \S\ref{sec:duasleis} Möbius aparece como
motivação --- é assim que o assunto se apresenta a quem chega. \textbf{A partir da prova das duas
leis, ele é consequência}: o que ali se demonstra não usa Möbius, e é Möbius que passa a sair dali.
\emph{Quem ler pela ordem do texto verá a mesma matriz duas vezes com estatutos diferentes, e é
deliberado.}

\medskip
\noindent As fracções contínuas são a realização onde tudo isto se mede.

\medskip
\noindent O texto está em \textbf{três partes que são duais entre si}, e o eixo é a dualidade de
Pontryagin ($\widehat\R=\R$, compacto\,$\leftrightarrow$\,discreto): a \textbf{álgebra} dá $+$,
$\times$ e inverso; a \textbf{topologia} dá ordem, limite e vizinhança; a \textbf{análise} dá
medida, taxa e distribuição. A tensão entre elas é medida e não decretada --- a álgebra opera sobre
conjuntos enumeráveis e não alcança a completude; a topologia alcança $\R$ e não fornece as
operações; a análise dá as leis de \emph{quase todo} número e por isso não vê nenhum dos objectos
aqui estudados, que são todos de medida nula.

\medskip
\noindent\textbf{O que este texto acrescenta, e é uma coisa só em três peças.}

\smallskip
\noindent\emph{(i) A definição:} a dualidade é a \textbf{memória da divisão}.
\emph{(ii) As duas leis}, bem tipadas: $1^{\dagger}=-1$ e $T^{\dagger}=-T$, com a marcação sem a
qual colapsam na solução trivial.
\emph{(iii) E o \textbf{corpo dual}}, que é o ambiente onde as leis têm significado:

\begin{center}\small
\begin{minipage}{0.92\textwidth}
\textbf{Definição.} \emph{O corpo dual é o par $(V,V^{*})$ juntamente com a identificação
$J\colon V\to V^{*}$.} As leis são enunciadas nele, \textbf{nunca num só dos seus lados} --- e é
por isso que não colapsam.
\end{minipage}
\end{center}

\noindent \textbf{É o corpo dual que unifica as três.} A definição diz \emph{o que} a dualidade
guarda; as leis dizem \emph{que forma} isso tem; e o corpo dual diz \emph{onde} isso pode sequer
ser dito. \emph{Sem ele as leis seriam igualdades mal tipadas --- ou triviais; com ele, são a
descrição de um objeto que tem dois lados e uma ponte.} E dentro dele prova-se que os
anti-autoadjuntos formam um corpo algébrico, autodual (\S\ref{thm:corpodual}) --- \emph{outra
palavra ``corpo'', outro sentido, e o texto distingue-as onde aparecem}. Da Segunda sai a \textbf{análise da família de
potência} (no \emph{Catálogo}): a escada é uma fórmula só e já continha a involução em $n=0$; o
coeficiente existe e é único para todo $n$, e a razão é uma desigualdade de inteiros; as duas raízes
\emph{são} o par $\sigma\sigma'=-1$; a paridade de $n$ decide quantas soluções reais existem; e
$(\sigma_n)_n$ fecha em $\Z[\sigma]$. Daí o teorema de que \textbf{a reversibilidade força Pisot em
grau dois} (no \emph{Catálogo}), com o critério lido nos coeficientes, $-A-1<B<A-1$. E a cadeia
$\N\to\dots\to\sigma$ passa a teorema, com as transições dimensionais a fecharem por bidualidade
(ambas no \emph{Catálogo}).

\smallskip
\noindent As \emph{equações} que a lei usa não são novas e estão citadas onde aparecem
(\cite{cook2021}; \cite{wonenburger1966,djokovic1967}): elas existiam separadas, e o que aqui se faz
é pô-las na mesma escada e analisar a família que delas resulta. As restantes atribuições vão
indicadas resultado a resultado, e vários resultados são deliberadamente \emph{negativos}: onde uma
analogia pára, é dito onde pára.
\subsubsection*{As operações da torre: clone e reprodução}

A torre não é só uma sequência de corpos: \emph{transporta estrutura de um nível para o
seguinte}, e as operações que a movem são duas.

\begin{center}\small
\begin{tabular}{@{}llll@{}}
\toprule
 & \textbf{os andares} & \textbf{a dimensão} & \\
\midrule
\textbf{clone}      & o \emph{mesmo}   & fica igual        & copia e \emph{não} gera \\
\textbf{reprodução} & \emph{diferentes} & vai a $\lcm(a,b)$ & gera e \emph{não} copia \\
\bottomrule
\end{tabular}
\end{center}

\noindent \emph{E a distinção é do próprio texto}, com as palavras que a fixaram: ``arquiteturas
diferentes é a ideia; se fosse igual seria clone''. O clone não move o tamanho --- é o lado
\textbf{aditivo}; a reprodução move-o pela $\lcm$ --- é o lado \textbf{multiplicativo}, e a conta
prova-o:
\[
  \lcm(a,b)\cdot\gcd(a,b) \;=\; a\cdot b .
\]
\textbf{A $\lcm$ é o produto com o comum descontado uma vez}, e $\gcd$ é o que os dois já tinham.

\medido{$\lcm\cdot\gcd=a\cdot b$ em $196$ pares, resíduo $0$ (\code{tests/bidual.c}).}

\paragraph{E o $\gcd$ que a $\lcm$ desconta é a memória da divisão.} Guardar só a $\lcm$ colapsa;
guardar o par não:

\medido{$144$ pares $(a,b)$ dão apenas $42$ valores de $\lcm$, mas $73$ pares
$(\lcm,\gcd)$ --- \textbf{$31$ distinções} que a $\lcm$ sozinha deitava fora. É a mesma contagem
que abre o texto, agora nas dimensões.}

\paragraph{E a dualização entre andares preserva a soma e espelha o produto.} A operação que liga
um andar ao seguinte é homomorfismo para $\oplus$ e \textbf{anti}-homomorfismo para $\otimes$:
\[
  \mathcal C(x\oplus y)=\mathcal C(x)\oplus\mathcal C(y),
  \qquad
  \mathcal C(x\otimes y)=\mathcal C(y)\otimes\mathcal C(x) .
\]
\emph{É a mesma inversão da adjunção, agora entre andares em vez de dentro de um} --- e é ela que
faz a torre alternar preto e branco.

\medido{$4096$ pares: a soma passa em todos, a ordem trocada passa em todos, e a ordem \emph{não}
trocada só em $472$ --- onde os dois comutam (\code{tests/bidual.c}).}

\noindent E as duas operações internas de cada andar já estão identificadas e medidas no
\emph{Catálogo}: $\oplus$ é Clifford e $\otimes$ é La Hire, com $36/36$ e resíduo $0$. \emph{O que
aqui se acrescenta não é a identificação --- é o lugar delas na torre.}

\subsubsection*{A forma tensorial: combinar corpos dá a terceira classe}

As operações do corpo cruz \emph{não são só cartesianas}: elas \textbf{combinam as classes dos dois
corpos e produzem uma terceira}. E a regra sai do produto dos discriminantes,
$\Delta(A\otimes B)=\Delta(A)\,\Delta(B)$ --- não é analogia, é a regra dos sinais:

\begin{center}\small
\begin{tabular}{@{}c|ccc@{}}
\toprule
$\otimes$ & hip. $(+)$ & par. $(0)$ & elíp. $(-)$ \\
\midrule
hip. $(+)$  & hiperbólico & parabólico & \textbf{elíptico} \\
par. $(0)$  & parabólico  & parabólico & parabólico \\
elíp. $(-)$ & \textbf{elíptico} & parabólico & \textbf{hiperbólico} \\
\bottomrule
\end{tabular}
\end{center}

\noindent \emph{Duas rotações compostas esticam}: $\text{elíptico}\otimes\text{elíptico}$ dá
\textbf{hiperbólico}, exatamente como $(-)(-)=(+)$. E o parabólico é o \textbf{zero absorvente}.

\paragraph{E para $n$ corpos, a lei é de uma linha.}
\[
  \Delta(A_1\otimes\cdots\otimes A_n) \;=\; \prod_i \Delta(A_i) ,
\]
donde: \emph{se algum fator é parabólico, o resultado é parabólico}; e senão, a classe decide-se
pela \textbf{paridade do número de elípticos} --- par dá hiperbólico, ímpar dá elíptico.
\textbf{É $(-1)^{\#\text{elípticos}}$}: o mesmo $(-1)^{n}$ de Cassini, do determinante e da torre.

\medido{$360$ combinações em $n=2\ldots5$, sem uma discordância com a lei
(\code{tests/bidual.c}; e a tábua de dois em \code{tests/viveiro\_metrico.c}).}

\paragraph{E a dimensão multiplica, com o logaritmo a somar.}
$\dim(A_1\otimes\cdots\otimes A_n)=\prod_i\dim A_i$, e $\log$ dessa dimensão é
$\sum_i\log\dim A_i$ --- \emph{o par aditivo/multiplicativo com $n$ fatores}. Nas potências de dois,
a soma dos $\log_2$ conta \textbf{quantos andares da torre se subiram}: $(2,3,5)$ dá dimensão $30$;
$(4,4,4)$ dá $64$, isto é seis andares.

\medskip
\noindent \emph{Uma nota de método, porque me apanhou:} ao medir isto pela primeira vez usei o
\emph{produto de matrizes} em vez do produto dos discriminantes, e as classes misturavam-se ---
$\text{hip}\otimes\text{elíp}$ dava hiperbólico em $80\%$ dos casos. \textbf{A operação era outra}, e
foi a tábua já medida no \emph{Catálogo} que o denunciou.


\part{Análise --- da métrica ao tempo}
\noindent\hfill\emph{\small[\textbf{realização} da Lei~2: o que se mede, e o que se move]}\medskip

\noindent Este capítulo constrói, pela ordem que a estrutura obriga, o que é preciso para
\emph{medir}: primeiro a métrica, depois o limite, depois a continuidade, e só então a derivada. E
termina onde a medida deixa de ser sobre um corpo e passa a ser sobre \emph{o conjunto dos corpos}
--- o corpo cruz e o corpo de corpos ---, que é onde a ordem da torre e a dinâmica do tempo se
encontram. \emph{É esse encontro que abre o capítulo seguinte.}

\subsubsection*{Ordenação e completude --- no corpo cruz, e não num lado só}

\emph{Dizer que $\C$ não é ordenável é verdade --- e é exatamente a metade que a dualidade existe
para não deixar sozinha.} A afirmação é sobre $\C$ \textbf{sem o seu par}; no corpo dual a pergunta
é outra, e a resposta também.

\paragraph{O corpo cruz, definido como estrutura e não por reutilização de símbolos.} Cada andar
não é $C_n$ nem $C_n^{*}$: é o par com a ponte \emph{e} as operações,
\[
  \boxed{\ \mathfrak C_n^{\times} \;:=\; \bigl(C_n,\ C_n^{*},\ J_n,\ \oplus,\ \otimes\bigr)\ }
\]
--- \emph{e escreve-se assim, e não ``$C+C^{*}$'' ou ``$C\times C^{*}$''}, porque esses símbolos já
significam soma direta e produto cartesiano e o objeto aqui é outro. \textbf{Nenhum dos lados é
completo isoladamente; a completude pertence ao par.}

\noindent E as duas \emph{operações do par sobre um elemento} --- que são o que faz a ordem existir
--- têm nomes próprios e caem nos reais:
\[
  z + z^{*} = 2\operatorname{Re}(z) \in\R \quad(\text{o \emph{traço}}),
  \qquad
  z \times z^{*} = |z|^{2} \in\R^{+} \quad(\text{a \emph{norma}}) .
\]
\emph{E reais ordenam-se.} \textbf{O corpo cruz é ordenado}, e a involução fá-lo naturalmente: ela
leva $z$ ao seu par, e o par cai onde há ordem.

\medido{Em $81$ inteiros de Gauss, a norma nunca é negativa; a ordem no par $(t,N)$ é total em
$2401$ comparações e compatível com a soma, sem uma falha (\code{tests/bidual.c}).}








\section{A construção dos corpos, e as operações que os movem}
\hfill \emph{\small[realização: exibe as duas coordenadas e mede o resíduo]}


\noindent\emph{Esta secção vem do texto narrativo, onde estava como apêndice técnico. É a
derivação corrida --- torção, assinatura, matrizes, transformada, convolução.}

\begin{tcolorbox}[colback=ouroclaro!35,colframe=ouro,boxrule=0pt,leftrule=2pt,arc=0pt,left=4mm,right=3mm,top=2mm,bottom=2mm,breakable]\itshape Vinte mundos atrás, oito aqui --- e a esta altura o leitor tem direito a desconfiar de que sejam vinte e oito coisas diferentes. \textbf{Não são.} Há \emph{um} corpo --- o universal, a raiz --- e há uma coisa só que gera todos os outros: a \textbf{assinatura}. Dita a assinatura, o corpo está determinado, e o veredito de se ele fecha lê-se nela. Como toda assinatura é uma \textbf{matriz}, cada corpo \textbf{é} uma matriz; reunidas, elas formam de novo um único corpo, o \textbf{Omnitrix}; e quem opera nele, sobre as matrizes, é o universal --- fazendo duas coisas, e só duas: a \textbf{transformada} (que vem em par dual: a \textbf{universal}, a fase, e a \textbf{dourada}, a taxa) e a \textbf{convolução}, que funde dois corpos num só. \textbf{Todas revertem} --- e é essa a razão de o conjunto fechar. \textbf{Depois disto não há mais prova: há realização.}\end{tcolorbox}

\subsection{Primeiro: o corpo universal --- a torção, e a base que não se escolhe}

O universal já está neste livro. Recolhe-se aqui o que o resto vai usar.

Escolhe-se um primo $p$, um $n$ que divide $p-1$, um gerador $g$, e a \textbf{torção}
\[
w=g^{(p-1)/n},
\]
o elemento cuja órbita fecha em $n$ passos. Toda a leitura vive em $\Z/p$: igualdade de inteiros, sem seno, sem cosseno, sem $\varepsilon$. Daí saem os \textbf{caracteres}, $\chi_k(j)=w^{jk}$, e eles não se escolhem --- não há base a inventar, porque \textbf{a base é a órbita da torção}. A única propriedade que se usa é
\[
\chi_k(u+v)=\chi_k(u)\,\chi_k(v):
\]
o caractere leva $\oplus$ a $\otimes$. É o $\prod$ de Pontryagin, e é tudo.

\paragraph{E a razão de a torção bastar.} Um elemento $w$ e a sua órbita chegam para mover o reino inteiro, e a razão cabe numa linha. Tome-se o dual de um grupo (os seus caracteres); tome-se o dual do dual. A aplicação natural
\[
\mathrm{ev}:\Gamma\longrightarrow\widehat{\widehat{\Gamma}},\qquad \mathrm{ev}_x(\chi)=\chi(x),
\]
--- avaliar $x$ em cada caractere --- é um \textbf{isomorfismo}: \textbf{o dual do dual devolve o grupo.} Não é resultado emprestado de fora: \emph{é a razão de a torção bastar}. Com o círculo, o dual são os inteiros, e o dual dos inteiros é o círculo outra vez: a reta que recobre e o círculo dos resíduos são um o dual do outro, e não há terceiro lugar para onde fugir.

Guarde-se esta frase, porque tudo o que se fizer daqui em diante volta a ela: \textbf{ir ao dual e voltar é a identidade.} Não são várias inversas avulsas --- a da transformada, a da convolução, a do flip. É \emph{uma} razão, e as inversas são consequências dela. É por isso que nada vaza.

Do lado do operador, o universal é o \textbf{sucessor}, $S(x)=x+1$: a soma é o sucessor iterado, a multiplicação é a soma iterada, e no bit o sucessor é o próprio complemento, $x\oplus1=\neg x$. Mas ele não se postula --- deduz-se. No princípio há apenas um instrumento que devolve, de cada objeto, um par $M(x)=(n,g)$ --- uma contagem inteira e um \textbf{resíduo de fechamento} --- e um caso distinguido, $g=0$. Nenhuma matriz, nenhuma operação suposta. E daí sai: \textbf{dois instrumentos com exatamente os mesmos fechamentos induzem necessariamente a mesma lei de composição.} A lei nasce das classes de equivalência dos fechamentos; o operador não foi inventado, é o único compatível com o princípio.

E uma distinção que faz falta adiante: as seleções válidas \textbf{não} formam o corpo --- formam a \emph{classe do neutro}. O corpo é o quociente \emph{inteiro}, dentro do qual o conjunto dos fechamentos é o elemento $1$. Quem tomasse só o quociente dos fechamentos obteria um grupo, pois todos os seus elementos têm resíduo nulo --- e jamais um corpo, onde $0\neq1$.

E a cascata que daí desce, porque a assinatura vai precisar dela: $\N$ (o monoide, de onde vem a ordem) $\to\Z$ (o anel, pelas classes de diferença) $\to\Q$ (o corpo, pelas classes de razão). Fica um degrau por subir --- de $\Q$ para a reta ---, e é a assinatura que o sobe.

\paragraph{E a cascata, degrau a degrau.} Ela desce toda do sucessor, e o que a faz
descer são as suas \emph{potências} --- quantas vezes se aplicou. Compor dá a soma,
$S^a\circ S^b=S^{a+b}$: os expoentes somam. Iterar dá o produto, $(S^a)^b=S^{ab}$: os
expoentes multiplicam. \textbf{As duas operações não são duas leis --- são a mesma
dinâmica lida de dois modos}, e é o $\exp/\log$ que troca uma pela outra.

A primeira torre é o sucessor sozinho: $0,1,2,\dots$ E com ela vem a \textbf{ordem}, sem
se acrescentar nada, porque o sucessor só avança --- $a\le b$ quando existe $c$ com
$a+c=b$. A segunda torre admite o antecessor, e dá o anel. Ali ainda não há inverso
multiplicativo: nenhum inteiro multiplica $2$ e dá $1$.

Os dois degraus seguintes são \textbf{classes de pares}, e a classe é o que apaga a
redundância. Pares com a mesma \emph{diferença} --- $(a,b)\sim(c,d)$ quando $a+d=b+c$ ---
dão os inteiros, com soma $(a+c,\,b+d)$ e inverso aditivo $(b,a)$, que é trocar as duas
camadas. Pares com a mesma \emph{razão} --- $(p,q)\sim(p',q')$ quando $pq'=p'q$ --- dão o
corpo, com produto $(pr,\,qs)$ e inverso multiplicativo $(q,p)$, que é trocar as duas
camadas outra vez. \textbf{O inverso é sempre a mesma troca; muda a camada em que ela se
faz.}

E fica a mesma falha do princípio: a classe está classificada e continua uma bagunça,
porque escolher um par dela é arbitrário. É por isso que o degrau seguinte não é mais uma
classe --- é a assinatura, que não escolhe.


Uma última coisa sobre o universal, e é a medida dele: a sua norma é $a^2$, e só. Uma coordenada, uma régua, e ela é \textbf{redonda}. \textbf{É a raiz.}

\subsection{Segundo: a ASSINATURA --- o que gera cada corpo}

Agora a peça de que sai o resto: \textbf{dita a assinatura, o corpo está determinado.} Não é uma etiqueta que se pendura num corpo já feito --- é o que o \emph{gera}. Quem tiver a assinatura tem o corpo, a régua, e o veredito de se aquilo fecha ou não fecha.

\textbf{A assinatura é a órbita.} Não é um rótulo que se pendura num corpo já feito, não é um quadro parado: é um \textbf{processo a correr}, e corre nos dois sentidos. Escreve-se como fração contínua,
\[
\frac ab=[a_0;a_1,\dots,a_n],
\]
que é medir ao resto zero --- tirar o maior número inteiro de vezes que cabe, e recomeçar com o que sobra.

Ela faz quatro coisas, e é preciso as quatro.

\textbf{Colapsa a classe.} Uma classe reúne infinitos pares, e escolher um deles é arbitrário --- a bagunça fica classificada, mas não gerada. A fração contínua não escolhe: ela some com o fator. A de $7/5$, a de $14/10$ e a de $21/15$ são a mesma, $[1;2,2]$ --- o $k$ desaparece no medir. É gerador, não representante.

\textbf{Gera-se por indução.} O passo é $x\mapsto a_i+1/x$, repetido; dele saem os convergentes, um a um, cada um mais perto que o anterior.

\textbf{No infinito, dá um metal.} Se o passo não termina, o limite não é racional: $[1;1,1,\dots]$ dá o ouro, $[2;2,2,\dots]$ a prata, $[3;3,3,\dots]$ o bronze. E o metal $\sigma_m=[m;m,m,\dots]$ é o autovalor dominante do gato $A_m$, de autovetor $[\sigma_m,1]$ e com $\sigma_m^2=m\sigma_m+1$ --- a direção que o gato não move, o eixo do boost. Os convergentes aproximam-se sem nunca chegar, e é assim que $\Q$ se completa: o degrau que faltava para a reta.

\textbf{E reverte.} Este é o lado que fecha o par, e sem ele a assinatura seria meia. Do metal, a mesma fração contínua \textbf{gera de volta} os convergentes --- $\sqrt2\to[1;2,2,\dots]\to \tfrac32,\tfrac75,\tfrac{17}{12},\dots$ --- isto é, gera a classe de onde veio. \textbf{Ida e volta.} A classe colapsa no metal; o metal regenera a classe. Quem calcula só a ida está a olhar meia folha.

E enquanto isso corre, os dois lados estão sempre abertos: um \textbf{expande} ($\sigma_m$), o outro \textbf{fecha} ($\sigma_m'=-1/\sigma_m$), e o produto deles é a norma, $\sigma_m\sigma_m'=-1$. A distância entre os dois, enquanto correm, é o \textbf{gap}, e ele sai da recusa do gato e do esquilo em trocar de ordem:
\[
\sqrt{m^2+4}\;=\;[\,\text{gato},\ \text{esquilo}\,].
\]
\textbf{A assinatura é o gap, e o gap é o comutador.}

\paragraph{E a conta é curta.} Multipliquem-se os dois nas duas ordens:
\[
A_mE=\begin{pmatrix}-1&m\\0&1\end{pmatrix},
\qquad
EA_m=\begin{pmatrix}1&0\\-m&-1\end{pmatrix},
\qquad
A_mE-EA_m=\begin{pmatrix}-2&m\\ m&2\end{pmatrix}.
\]
A diferença tem traço zero e determinante $-(m^2+4)$, logo as suas raízes são
$\pm\sqrt{m^2+4}$ --- que é a distância entre as duas raízes de $x^2-mx-1$. E o
anticomutador dá $A_mE+EA_m=mE$: a ordem do metal, pura. \textbf{A recusa em trocar de
ordem é a assinatura; a concordância é a ordem.}


\paragraph{E a fotografia dela.} Parada a órbita num instante, o que fica é a norma --- $N(x)=x^{\!\top}\!Ax$, com $A$ simétrica --- e medir passa a ser uma conta de matriz. A do universal é a de uma entrada só, $A=(1)$, que dá $N=a^2$. É a raiz.

O universal impõe uma parte e deixa outra por preencher, e a construção mora nessa fresta. O corpo universal não é uma família numérica nova: é a \textbf{estrutura mínima} imposta pela exigência de que um instrumento feche exatamente sobre o objeto que mede. Toda linguagem que satisfaça esse princípio é uma \emph{realização particular} da mesma estrutura --- e o que o princípio impõe, sozinho, é a \textbf{lei de composição}. \emph{A soma é o que as realizações acrescentam.} Daí a pergunta que a assinatura responde: \textbf{o que foi que esta realização acrescentou à raiz?}

E toda matriz simétrica é, a menos de mudança de base, uma soma de cópias dessa régua da raiz, com sinal. Diagonaliza-se, e o que sobra na diagonal são $+1$, $-1$ e $0$ --- a régua $a^2$ repetida, ora somando, ora subtraindo, ora calada. Contem-se os dois lados, e o que colapsou:
\[
(p,\,q,\,r)\;=\;\bigl(\text{quantas réguas redondas},\ \text{quantas hiperbólicas},\ \text{quantas colapsadas}\bigr),
\]
e o grau é $p+q+r$. \textbf{Um corpo é a raiz repetida $p+q+r$ vezes, cada cópia com o seu sinal} --- e a contagem diz, número a número, o que aquela realização acrescentou ao $(1,0,0)$ de onde todas partem. É leitura, não definição: a órbita continua a correr por baixo, e é dela que estes números são o retrato. O reflexivo é $a^2+c^2$, assinatura $(2,0,0)$: duas redondas. O expansivo é $a^2-Db^2$ com $D>0$, assinatura $(1,1,0)$: uma redonda e uma hiperbólica. O geométrico é $t^2-x^2-y^2-z^2$, $(1,3,0)$. O eletromagnético é $E^2-B^2$, $(3,3,0)$. O telescópico é $(2,2,0)$; o conforme, $(4,1,0)$, grau cinco; o celeste, $r^2+C^2=1$, $(2,0,0)$; o óptico, $(2,0,0)$; o cristalino, $a^2+|D|b^2$ com $D<0$, $(2,0,0)$; o relógio, $1-s^2$, $(1,1,0)$; o sensitivo e o evolutivo, $(1,1,0)$ os dois; o individual, $(2,0,0)$; o deflexivo, $(1,0,0)$ --- a mesma da raiz. Nenhum destes é um mundo à parte, e nenhum destes números é uma etiqueta: são leituras de dinâmicas diferentes, a mesma régua tomada tantas vezes e com tantos sinais.

E a fotografia não perde o que a norma sabe: dois corpos com a mesma contagem são \textbf{congruentes} --- isto é, se um é o outro visto noutra base: qualquer troca de coordenadas $P$ invertível leva $A$ em $P^{\!\top}\!AP$ e \textbf{preserva} $(p,q,r)$. Nada mais se extrai da norma além desses três números --- o que ela sabe, sabe-o ali; o que ela não vê é o movimento de onde veio.

\subsubsection{O critério --- quando uma matriz é corpo, e quando não é}

Falta o teste. Uma matriz $2\times2$ tem um \textbf{dual} que vive dentro da própria álgebra, porque é combinação da matriz com a identidade:
\[
\bar M=\operatorname{tr}(M)\,I-M,\qquad\text{e daí}\qquad M^{-1}=\frac{\bar M}{\det M}.
\]
O inverso é o dual dividido pela norma. Portanto tudo depende de uma coisa só --- \textbf{se a norma se anula ou não}:

\begin{tcolorbox}[colback=ouroclaro!35,colframe=ouro,boxrule=0pt,leftrule=2pt,arc=0pt,left=4mm,right=3mm,top=2mm,bottom=2mm,breakable]\itshape Uma álgebra de matrizes é um \textbf{corpo} se, e somente se, a sua norma é \textbf{anisotrópica} --- isto é, se $N(x)=0$ só acontece em $x=0$.\end{tcolorbox}

Se algum elemento não nulo tem norma zero, ele não inverte, e acabou. Se a norma nunca zera, o inverso já está lá dentro, pela fórmula acima --- não é preciso ir buscá-lo fora.

E o teste diz mais: diz \emph{de que tipo} é cada álgebra. O determinante, como forma, tem assinatura $(2,2)$ --- dois sinais de cada lado ---, e por isso nenhuma parte anisotrópica dele passa do grau $2$. Daí os corpos serem pequenos. Dentro do grau $2$, quem decide é o \textbf{discriminante} $\Delta=\operatorname{tr}(M)^2-4\det M$, e há exatamente três desfechos:

\begin{itemize}
\item $\Delta<0$ --- o regime \textbf{elíptico}, o da torção: o grupo é o \textbf{círculo}, e a álgebra \textbf{é corpo}. É o único dos três que fecha, e o único cujo movimento a um parâmetro é compacto: a órbita volta.
\item $\Delta=0$ --- o regime \textbf{nilpotente}: o grupo é a \textbf{reta}, há um elemento não nulo cujo quadrado é zero, e não há inverso. É o \textbf{glacial} --- e é exatamente onde mora o corpo exterior, com o seu $e^2=0$.
\item $\Delta>0$ --- o regime \textbf{idempotente}: o grupo é a \textbf{hipérbole}, e a álgebra parte-se em duas retas independentes. É o \textbf{tropical} --- o mundo do $\max$, onde nada se desfaz.
\end{itemize}

\textbf{Apenas a torção fecha um corpo.} Sobre os racionais, quem impede a norma de zerar é a irracionalidade do metal --- a assinatura que não termina. \textbf{Um corpo é uma matriz cuja norma não sabe chegar a zero.}

E daí saem três coisas.

\paragraph{Quando o corpo morre.} $r>0$ significa que alguma direção colapsou --- há um elemento não nulo de norma zero, portanto um divisor de zero, portanto não há inverso. \textbf{$r>0$ $\iff$ a forma degenera $\iff$ o corpo morre.} De todo o bestiário, um só tem $r=1$: o \textbf{exterior} ($N(a+be)=a^2$, com $e^2=0$). Todos os outros têm $r=0$ --- o que não os torna corpos a todos: o \textbf{conforme} tem $r=0$ e grau cinco, $(4,1,0)$, e mesmo assim não fecha, porque com sinais dos dois lados a norma encontra o zero fora do zero. É o \textbf{cone nulo}, e é por isso que o conforme é a \emph{casa} da medida e não um corpo. É por isso que o exterior, sozinho, não é corpo --- e por isso o nome dele nomeia um degrau, não um mundo acabado.

\paragraph{Estar na lista não basta.} O grau tem de ser $1$, $2$, $4$ ou $8$ --- os degraus da torre, que termina nos octoniões, de assinatura $(8,0,0)$: oito réguas redondas, o corpo final, e oito é o período de tudo neste livro. Mas a condição é \textbf{necessária, não suficiente}: o exterior tem grau $2$, está na lista, e mesmo assim não fecha, porque lhe falta a anisotropia. Estar na lista é permissão de entrar; o que decide é o \textbf{sinal} da norma --- radical, não; definida, sim.

\paragraph{E há corpos onde a fotografia não sai.} A contagem $(p,q,r)$ precisa de uma norma de grau dois para a que aplicar, e nem todo corpo tem uma. O \textbf{mórfico} mede pelo peso de contagem sobre $\mathbb{F}_2$: ali a homogeneidade $q(\lambda x)=\lambda^2q(x)$ não vale. O \textbf{espaço-temporal} tem régua \textbf{ultramétrica}, $|x+y|\le\max(|x|,|y|)$ --- na régua comum ter-se-ia $|1+1|=2>\max(1,1)$, e a desigualdade forte cai. O \textbf{entrópico} tem $\otimes=+$: \emph{a multiplicação é a soma}, e não sobra norma. O \textbf{local} e o \textbf{global} medem por divisibilidade e pela fórmula do produto, não sobre a reta. E o Rei: a régua dele é de \textbf{grau 4}, $\mathrm{sl}^2+\mathrm{cl}^2+\mathrm{sl}^2\mathrm{cl}^2=1$, e o termo cruzado $\mathrm{sl}^2\mathrm{cl}^2$ não tem sinal onde ser contado --- \textbf{o Rei está acima da régua.}

Nada disto os põe fora da construção; põe fora da \emph{fotografia}. A órbita continua a correr em todos eles --- e onde a contagem não alcança, alcança o dual: \textbf{nenhum corpo opera sozinho.} O caso está escrito neste livro: o entrópico, que sozinho não tem inverso aditivo, recupera-o sob o exponencial do tabuleiro cósmico. E repare-se que é a mesma coisa duas vezes --- o que lhe tira a norma ($\otimes=+$, o produto colapsado na soma) é exatamente o que o logaritmo fez ao produto. O que se perdeu de um lado, o outro devolve.

É esta a ligação que falta dizer, e ela é dupla. \textbf{O dual} é o parceiro que dá o que falta; \textbf{a transformada} é o caminho que leva de um ao outro. Um corpo isolado não opera --- não por fraqueza, mas porque a conservação nunca viveu dentro de um corpo: vive sobre o par. É isso que faz do conjunto um corpo e não uma prateleira: os corpos não estão lado a lado, estão \textbf{ligados}, e ligados de forma que se pode ir e voltar.

\subsection{Terceiro: o Omnitrix --- o corpo que os corpos formam}
\hfill \emph{\small[realização da Lei~2 (a dualidade é dual): o conjunto fecha porque todas revertem]}


Reunidos, os corpos não fazem uma coleção: fazem um corpo. É a esse que se dá o nome de \textbf{Omnitrix}, e a definição tem de ser dita com cuidado, porque a palavra promete menos do que a coisa é:

\begin{tcolorbox}[colback=ouroclaro!35,colframe=ouro,boxrule=0pt,leftrule=2pt,arc=0pt,left=4mm,right=3mm,top=2mm,bottom=2mm,breakable]\itshape \textbf{O Omnitrix não é um tabuleiro: é o \emph{fractal} de tabuleiros que o axioma mantém coeso.}\end{tcolorbox}

Não é a caixa onde os corpos são guardados nem o índice deles. É o que sobra quando se percebe que cada corpo, olhado de perto, é feito de corpos --- e que a lei é a mesma em toda escala. Os reais, os complexos, os quatérnios, os octoniões e os metais entram nele como degraus, com \textbf{duas} operações e um invariante:

\begin{itemize}
\item $\oplus$ é a \textbf{soma direta}: pôr um corpo ao lado do outro. O resultado tem duas partes, e o grau é a soma dos graus.
\item $\otimes$ é \textbf{La Hire}: o produto, o rolar, aquilo que compõe.
\item os geradores são \textbf{dois}, e apenas dois: o \textbf{gato} $A=\bigl(\begin{smallmatrix}1&1\\1&0\end{smallmatrix}\bigr)$, simétrico, $\det=-1$, com $A^2=A+I$, autovalores reais, que \emph{dissipa}; e o \textbf{esquilo} $E=\bigl(\begin{smallmatrix}0&1\\-1&0\end{smallmatrix}\bigr)$, antissimétrico, $\det=+1$, com $E^2=-I$, autovalores no círculo, que \emph{conserva} --- é a torção, e é o seu período $4$ que já se viu na transformada universal, $\F^2=$ reflexão, $\F^4=\mathrm{id}$.
\item o \textbf{invariante} é a torção; e a \textbf{lei de conservação} é a norma:
\[
\det\,(X\otimes Y)=\det X\cdot\det Y .
\]
A norma de um produto é o produto das normas. Uma linha --- e é ela que proíbe o vazamento em todo o Omnitrix.

\paragraph{E por que dois bastam.} O gato aplicado a um par é o meio-somador,
\[
A_1(x,y)=(x+y,\ x),
\]
e a componente de cima tem \textbf{duas} leituras, que são as duas operações: o que sobra
ao dividir por dois é a soma, e o que passa é o produto,
\[
(x+y)\bmod 2 \;=\; \oplus,
\qquad
(x+y)\operatorname{div} 2 \;=\; \otimes .
\]
Não se acrescentaram duas operações ao gato. \textbf{São as duas projeções dele.}

E delas sai o resto sem primitiva nova. A negação é $x\oplus1$. A disjunção sai sem sequer
a negação,
\[
a\vee b=a\oplus b\oplus(a\wedge b),
\]
que se lê nos quatro pontos --- em $(1,1)$ dá $1\oplus1\oplus1=1$. Daqui vem a lei do
núcleo do reino, e ela não é voto de pobreza: \textbf{a divisão não é proibida, é
supérflua} --- o gato já entrega as duas operações que geram tudo o que é finito. Quem
escreve um núcleo com divisão está a consumir a operação que o núcleo devia gerar.

\paragraph{E a conservação, em inteiros.} A potência do gato tem as entradas do Fibonacci
de ordem $m$,
\[
A_m^{\,k}=\begin{pmatrix}F_{k+1}&F_k\\F_k&F_{k-1}\end{pmatrix},
\qquad
\det\bigl(A_m^{\,k}\bigr)=(\det A_m)^k=(-1)^k,
\]
e lendo o determinante nas entradas, $F_{k+1}F_{k-1}-F_k^2=(-1)^k$. A mesma quantidade
lida como norma: de $\sigma_m^{\,k}=F_k\sigma_m+F_{k-1}$ sai
$F_{k-1}^2+m\,F_{k-1}F_k-F_k^2=(-1)^k$. O par é $(F_{k-1},F_k)$, a segunda coluna;
trocá-lo pela primeira quebra a igualdade. \textbf{A lei de conservação do Omnitrix, aqui,
é uma igualdade de inteiros --- sem tolerância nenhuma.}

\end{itemize}

\paragraph{O que faz dele um fractal, e não uma pilha.} Três peças, e elas se sustentam umas às outras.

Primeiro, a \textbf{promoção é o idempotente}: $x^2=x$. Quem chega ao fim promove-se, e promover-se duas vezes é promover-se uma --- o cristal, a única classe de ponto fixo. Segundo, \textbf{ser fail-closed é a mesma coisa que promover.} Num tabuleiro finito e sozinho, o peão que atinge a borda não tem casa seguinte: acabou, e essa é a porta fechada. Mas o tabuleiro não está sozinho: o que fecha num é a promoção no outro, e a norma composta continua a valer, $N(B_1\otimes B_2)=N(B_1)\,N(B_2)$. \textbf{Um fail-closed num tabuleiro é redimido por uma promoção no outro} --- e é por isso que a borda não é o fim do mundo, é uma costura.

Terceiro, e é o que fecha: a lei atravessa \textbf{toda} escala. Tome-se por hipótese que o axioma da norma vale no nível $k$, com a promoção $x^2=x$ e a redenção acima. O nível seguinte toma os tabuleiros do nível $k$ como \emph{peças}; e como a dinâmica é a mesma por unidade, a subida de um nível ao outro é uma \textbf{homotetia} --- uma mudança de escala pura. Ora, a parábola é invariante sob homotetia. Logo o axioma atravessa, e atravessa para sempre. É uma indução cujos objetos são, eles próprios, induções: uma \textbf{indução de induções}. E a conclusão é a definição: o fractal de tabuleiros é \textbf{fechado sob a parábola em toda escala}.

Um axioma, uma torre, um reticulado, uma partida, um tempo --- e a mesma estrutura em toda escala. \textbf{É isso o Omnitrix.}

E é aqui que o gap de há pouco --- a assinatura, o comutador do gato com o esquilo --- se mostra como a coisa que roda igual em toda escala: solte-se a dimensão, solte-se o substrato, e o que sobra é ele. O fractal não é feito de tabuleiros parecidos: é feito da \emph{mesma} recusa, repetida.

\subsection{Quarto: as matrizes --- cada assinatura gera o seu corpo}

Agora a assinatura a produzir o corpo, e é uma conta só. Tome-se a assinatura $m$. Ela dá o gato $A_m=\bigl(\begin{smallmatrix}m&1\\1&0\end{smallmatrix}\bigr)$, cujo autovalor $\sigma_m$ satisfaz a equação do gato, $\sigma_m^2=m\,\sigma_m+1$. E então
\[
K_m=\{\,a+b\,\sigma_m \;:\; a,b\in\Q\,\}
\]
\textbf{é um corpo}. Some-se, multiplique-se, inverta-se: fecha nos três.

A soma é componente a componente. O produto fecha porque a equação do gato reescreve o quadrado:
\[
(a+b\sigma)(c+d\sigma)=ac+(ad+bc)\sigma+bd\,\sigma^2=\underbrace{(ac+bd)}_{\in\Q}+\underbrace{(ad+bc+m\,bd)}_{\in\Q}\,\sigma .
\]
\textbf{É a única regra do corpo, e é o gato.} Nenhum grau novo aparece: o $\sigma^2$ desce sempre a grau $1$. Comutar, associar, distribuir, achar o neutro --- tudo isso vem de graça dos racionais, porque a soma é a deles e a redução do quadrado é sempre a mesma.

Falta o que separa um anel de um corpo: \textbf{todo elemento não nulo inverte.} Quem o dá é o esquilo. A raiz conjugada satisfaz a mesma equação, e dela saem duas relações:
\[
\sigma+\sigma'=m,\qquad \sigma\,\sigma'=-1 .
\]
Defina-se o conjugado $\bar u=a+b\sigma'$ --- \textbf{a troca de sinal fora da diagonal, que é o esquilo}. Então a norma é racional, e o inverso está escrito:
\[
N(u)=u\,\bar u=a^2+m\,ab-b^2,\qquad u^{-1}=\frac{\bar u}{N(u)} .
\]
E aqui as duas assinaturas deste livro encontram-se. Para o inverso existir, é preciso que $N(u)$ nunca se anule fora do zero. Ora, $N(u)=0$ com $u\neq0$ exigiria que $\sigma_m$ fosse racional --- e $\sigma_m$ é o limite da assinatura $[m;m,m,\dots]$, \textbf{que não termina}. \emph{É a assinatura da órbita que garante que o corpo fecha.} A fração contínua não é ornamento da construção do número: é a razão pela qual o corpo tem inverso.

E nada disto usou o valor de $m$ --- usou-se apenas $\sigma^2=m\sigma+1$. Logo há um corpo $K_m$ \textbf{para cada $m$}, todos pela mesma conta. \textbf{A construção é uma só; a assinatura escolhe a classe.}

\subsection{A primeira operação: a TRANSFORMADA UNIVERSAL --- e a sua inversa}

Construídos os corpos, falta o que se faz com eles. E o que o universal faz, agindo no Omnitrix sobre as matrizes, são duas operações. A primeira leva um corpo a \textbf{outro}, e chama-se \textbf{transformada universal} porque é do corpo universal que ela é feita --- não há nada nela que venha de fora: os caracteres são a órbita da torção, e a torção é dele.

E aqui está a coisa que o leitor deve levar antes da fórmula, porque é ela que explica por que o universal \emph{opera} sobre as matrizes em vez de morar ao lado delas: \textbf{a transformada não se define --- ela cai do dual.} Já se viu que o dual do dual devolve o grupo; então o emparelhamento de um objeto com um caractere não é uma construção a mais, é a \emph{avaliação} dos caracteres --- e os caracteres são a órbita da torção. Daí a consequência que importa: \textbf{trocar o grupo troca a transformada sem trocar a conta.} O que o corpo universal fornece não é uma transformada: é o \textbf{mecanismo que as produz todas}. Ele não é mais uma matriz ao lado das outras --- é quem gera a transformada de cada uma.

Vai assim, e volta assim:
\[
\F(x)_k=\frac{1}{\sqrt n}\sum_j x_j\,\chi_k(j),
\qquad
\F^{-1}(X)_j=\frac{1}{\sqrt n}\sum_k X_k\,\chi_{-k}(j).
\]
A normalização é $1/\sqrt n$ em \emph{cada} lado, e não por gosto: somar a órbita inteira dá $\sum_k\chi_k(j)\chi_{-k}(j')=n\,\delta_{j,j'}$ --- o ponto ---, e esse $n$ tem de ser cancelado por duas metades, uma de cada aplicação. É a \textbf{ida e a volta}, e a volta é exata: nenhuma das duas é aproximação de coisa nenhuma.

E há mais do que a inversa: a transformada universal é \textbf{periódica}. Aplicada quatro vezes devolve o que entrou; aplicada duas, devolve a reflexão:
\[
\F^2=\text{reflexão},\qquad \F^4=\mathrm{id}.
\]
É o \textbf{mesmo período $4$ do esquilo} --- a torção que gira, tem período e não dissipa. Três batidas de um lado são a volta. Por isso os corpos não são vizinhos numa lista: são transformadas uns dos outros, e o que separa dois deles é o número de aplicações entre um e outro. Quem sabe a transformada universal não decora a tabela: vai de um ao outro, e volta quando quiser.

E há uma razão dura para a volta ser exata, e não aproximada: \textbf{a transformada é unitária} --- ela preserva a norma,
\[
\|x\|^2=\|\F x\|^2 ,
\]
e quem preserva a norma tem inversa, e a inversa não perde nada. Dito em português do reino: medir um objeto é decompô-lo em pedaços da unidade, e aqui os pedaços são os caracteres, isto é, as iteradas da torção; a identidade acima diz que a soma dos pedaços reconstitui o objeto \textbf{sem sobra e sem excesso}. \textbf{Nenhuma medida vaza na decomposição.} Não é coisa nova: é o Princípio da Medida Consistente lido nos coeficientes --- a mesma conservação da unidade, agora na base da torção. E ela reaparece em toda parte com outra roupa: a torção tem determinante $+1$, o Rei ocupa sempre exatamente uma casa, e numa álgebra normada a forma multiplicativa da mesma frase é $\|xy\|=\|x\|\,\|y\|$.

\subsubsection{A outra metade: a TRANSFORMADA DOURADA --- a taxa}

O corpo universal tem \textbf{duas} metades, e cada uma tem a sua transformada. Não é uma escolha de método: é a decomposição do próprio grupo, que se parte num fator de \textbf{taxa} e num fator de \textbf{fase}.

\begin{itemize}
\item \textbf{o eixo angular --- a fase.} O grupo é o círculo, a medida invariante é $d\theta/2\pi$, o dual são os inteiros, e a transformada é a \textbf{universal} --- a que se acabou de escrever, a que pergunta ao mundo \emph{de que ondas ele é feito}.
\item \textbf{o eixo radial --- a taxa.} O grupo é a reta multiplicativa, a medida invariante é $dt/t$, o dual é a reta, e a transformada é a \textbf{dourada} --- a que pergunta ao mundo \emph{de que tamanho ele é}.
\end{itemize}

E a segunda não é uma segunda máquina: é a primeira, \textbf{flipada}. O \textbf{flip} é a mudança de reta --- escreve-se $u=\log t$, e a reta onde se multiplica passa a ser a reta onde se soma. A medida vai junto: $dt/t$ vira $du$. Aplique-se então a transformada que já se tem, do outro lado da mudança:
\[
\F\circ\log ,
\]
e o que sai é \textbf{a transformada dourada} --- unitária por ser composição de unitárias, e por isso também ela inverte sem perda. \textbf{O flip é que a dá.} Não são duas transformadas: é uma, e a outra reta. É a régua de Merlin, a de base $\varphi$: a lente girassol, que é invariante por escala áurea --- muda-se o \emph{zoom} e a resposta é a mesma resposta.

E a inversa dela sai da mesma maneira: a inversa de uma composição é a composição das inversas, na ordem trocada. Escrevendo $\Lambda$ para o flip, $(\Lambda x)(u)=x(e^u)$, com $(\Lambda^{-1}y)(t)=y(\ln t)$:
\[
\mathfrak{G}=\F\circ\Lambda,\qquad \mathfrak{G}^{-1}=\Lambda^{-1}\circ\F^{-1} .
\]
Desfaz-se o giro, depois desfaz-se a escala. Nada se perde no caminho, porque as duas metades são unitárias e a unitariedade compõe.

O livro já tinha dito isto na colina, e agora tem a fórmula: \textbf{quem pusesse o mundo da frequência em escala logarítmica, $x=e^u$, veria a régua de um virar, ponto por ponto, sem sobra e sem resto, a régua do outro.} São a mesma transformada em coordenadas diferentes --- o sábio da escala e o adversário das ondas, um par, e nenhum dos dois a saber.

E agora o par pode ser dito com precisão, porque é o par que abre o livro. \textbf{Trocar o grupo troca a transformada sem trocar a conta}: sobre o grupo \textbf{aditivo} a transformada é a das ondas --- o lado $\oplus$, a frequência, o adversário da planície; sobre o grupo \textbf{multiplicativo} é a da escala --- o lado $\otimes$, a régua do sábio da colina; \textbf{e a segunda é a primeira composta com o logaritmo}. Não são dois métodos rivais: são as duas metades duais de um mesmo corpo, e a passagem entre elas é $\prod=\exp\circ\Sigma\circ\log$, a costura que este livro vem repetindo desde o Batimento Único. O elenco não estava pintado por cima da matemática --- \textbf{os dois adversários são as duas metades.}

\paragraph{E o processo segue.} O flip leva de uma reta à outra; a transformada trabalha ali; o flip leva de volta. Nada nesse vaivém encontra uma parede --- \textbf{o processo não termina, segue.}

E essa não-terminação já apareceu aqui, com outra roupa. A assinatura de um metal é
\[
\sigma_m=[m;m,m,\dots],
\]
que é o passo $x\mapsto m+1/x$ repetido para sempre --- e dentro desse passo está o flip, a inversão $1/x$. O metal é o \textbf{ponto fixo} dele: a cifra que a repetição não move. Se a repetição parasse, o metal seria racional; e se fosse racional, a norma encontraria o zero fora do zero, e o corpo não fecharia.

Ou seja, é a mesma coisa dita duas vezes: \textbf{o processo não terminar é o que faz o corpo fechar.} A infinitude não é o preço do fechamento --- é a condição dele.

Portanto: \textbf{taxa $\bowtie$ fase, dourada $\bowtie$ universal.} Cada metade tem a sua conservação, e o produto das duas \emph{é} a dualidade de Pontryagin --- o mesmo par que já se viu no motor e no rotor, e que se escreve numa linha, $\log z=\log|z|+i\theta$: a magnitude, que é aditiva, e a fase, que gira. Um corpo do bestiário pode ser lido pela frequência ou pela escala; as duas leituras revertem, e as duas dão o mesmo corpo.

\subsection{A segunda operação: a CONVOLUÇÃO --- e a sua inversa}

A segunda operação \textbf{combina} dois corpos. Dados $a$ e $b$,
\[
(a\conv b)_j=\sum_{u+v\equiv j}a_u b_v,
\qquad\text{e}\qquad
\F(a\conv b)=\sqrt n\;\F(a)\cdot\F(b).
\]
Sai numa linha, e a linha é a única propriedade que o universal tem. Escreva-se a transformada da convolução e troque-se a ordem das somas: aparece $\chi_k(u+v)$, que é $\chi_k(u)\chi_k(v)$, e a soma dupla \textbf{separa em duas somas independentes} --- uma em $u$, outra em $v$. \emph{A soma dupla separa porque o caractere separa.} Nada mais entra na conta; não há aproximação, e a igualdade é de inteiros.

E daqui sai o desfazer: \textbf{do outro lado, a convolução é um produto ponto a ponto.} Do lado do ponto, fundir custa todos os pares; do lado do caractere, custa uma multiplicação por casa. \textbf{Fundir é multiplicar do outro lado --- e por isso desfazer é dividir.} O mesmo passo aparece na luz: pôr uma máscara sobre o objeto multiplica o campo, e multiplicar no objeto é convolver no espectro. Produto e convolução são as duas faces de um lance só.

E repare-se no que isto diz. \textbf{Fundir não é pôr um ao lado do outro} --- isso é $\oplus$, a soma direta, e o resultado teria duas partes: o grau cresce e as assinaturas somam-se, que é como \emph{dois corpos geram um terceiro}, pois a assinatura de uma soma é a soma das assinaturas. Fundir é $\otimes$: \textbf{todo ponto de um encontra todo ponto do outro}, e o que sai tem o tamanho de um só. São duas maneiras de combinar, e a diferença entre elas está no grau: ao lado, o grau soma; fundidos, o grau fica.

Para $L$ corpos, um só produto:
\[
\F(a_1\conv\cdots\conv a_L)=(\sqrt n)^{L-1}\prod_{l=1}^{L}\F(a_l).
\]
$L$ corpos, $L-1$ costuras, uma casca.

\textbf{E desfaz-se.} Se nenhum dos outros zerou,
\[
a_m=\F^{-1}\!\left(\frac{\F(C)}{(\sqrt n)^{L-1}\prod_{l\neq m}\F(a_l)}\right)
\]
exato, sem resíduo. \textbf{Quem abre é o inverso}, e o inverso é o degrau que leva do anel ao corpo --- é ele que faltava. Sem ele haveria como fechar e não haveria como abrir, e o que a corte teria construído não seria uma transformação: seria um cofre.

\paragraph{E o único lugar onde não há volta.} É o zero, e a corte já esbarrou nele no alto da torre: uma casa transformada que dá zero fica \textbf{selada}, porque o zero é o único elemento sem inverso. O que a derivação acrescenta ao que lá se viu é o alcance do estrago --- ele \emph{parece} local e não é. Ponha-se $\varepsilon$ no lugar do zero e olhe-se o que volta:
\[
\F^{-1}(Q)_j-a_j=\frac{1}{\sqrt n}\,\varepsilon\,w^{-kj},
\]
que não se anula para nenhum $j$: o erro de uma casa espalha-se por todas, com a mesma amplitude, e nenhuma escapa. \textbf{Uma casa selada não estraga uma letra: estraga o texto inteiro.} No corpo não existe ``quase certo'' --- volta exato, ou volta todo errado.

\subsection{As realizações --- os oito corpos que completam o bestiário}

E agora o mais curto. \textbf{Daqui para a frente não há prova nenhuma} --- a prova foi uma, e foi a construção: a raiz, a assinatura, o critério. Cada um dos oito que seguem, e cada um dos vinte que vieram antes, é \textbf{a mesma coisa noutro idioma}. O que se escreve sobre cada um não é justificação: é \textbf{realização} --- a sua assinatura, o que ela lhe dá e o que lhe tira, e como aquilo se joga.

É também por isto que ninguém entra de fora. Depois da construção não sobra porta: não há passo à espera de ser importado, porque não há passo que não tenha sido derivado.

\subsubsection{o tabuleiro --- as peças são o Rei}

Aqui a construção joga-se, e a régua é a do peso, sobre $\mathbb{F}_2$: é dos que não têm
fotografia --- $(p,q,r)$ não se aplica, e o que mede é a contagem.

Há \textbf{uma} peça, e é a vizinhança de um passo:
\[
\mathcal{D}=\{(\pm1,0),\ (0,\pm1),\ (\pm1,\pm1)\},
\]
oito direções, que módulo o sinal são quatro retas. Dela saem as duas operações, e elas
correm em sentidos contrários. \textbf{Dilatar} espalha,
$\delta_B(A)=\bigcup_{b\in B}(A+b)$; \textbf{erodir} retém só o que sobrevive a todos os
deslocamentos ao mesmo tempo, $\varepsilon_B(A)=\{x:\ x+b\in A\ \text{para todo}\ b\in B\}$.

Uma peça é uma \textbf{máscara} e uma \textbf{potência}. Com a máscara cheia e potência
um, o Rei; com ela e potência sem limite, a Dama; com as duas retas ortogonais, a Torre;
com as duas diagonais, o Bispo. Nada as distingue senão a potência. E as retas da Torre e
do Bispo não se cruzam --- uma casa ortogonal partilha exatamente uma coordenada com a
origem, uma diagonal tem as duas diferenças iguais e não nulas, e nenhuma casa é as duas
coisas ---, o que faz da Dama a soma \emph{exclusiva} das duas: das $56$ casas que ela
alcança de um canto livre, $28$ são de uma e $28$ da outra.

\paragraph{E não existe Cavalo.} Iterar a dilatação dá quadrados: $\delta^2$ é o
$5\times5$, $\delta^1$ é o $3\times3$, e a diferença tem $25-9=16$ casas. Dessas
dezasseis, oito são múltiplos por dois de uma das oito direções --- $(\pm2,0)$,
$(0,\pm2)$, $(\pm2,\pm2)$ --- e uma peça de reta já lá chega. Sobram oito,
$(\pm2,\pm1)$ e $(\pm1,\pm2)$, com as duas coordenadas não nulas e diferentes em módulo,
logo múltiplos de direção nenhuma. \textbf{Essas oito são os saltos do Cavalo.} Ele não é
peça nova nem exceção: é o que a segunda potência do Rei alcança e nenhuma reta alcança.

\paragraph{E só o Peão tem seta.} A sua máscara de captura, $\{(1,1),(1,-1)\}$, não contém
as antípodas --- é a única que distingue sentido de direção. Daí a promoção: para quem é a
sua própria antípoda, ir e voltar é a mesma reta, e a oitava fileira não existe como
destino. Daí também a cor, porque a antípoda da máscara branca é a preta. E é dele que o
Rei se levanta: juntando a máscara à sua antípoda saem três retas, e a quarta é a
composição dos dois fluxos do próprio Peão --- $\{(-1,0)\}$ somado a $\{(1,1),(1,-1)\}$ dá
$\{(0,1),(0,-1)\}$.

\paragraph{E as duas operações revertem.} É a lei deste capítulo, escrita em máscaras:
erodir por uma é dilatar pela \textbf{virada}, no complemento,
\[
\varepsilon_B(A)=\neg\,\delta_{\check{B}}(\neg A).
\]
Um ponto está na erosão quando não existe deslocamento que o tire de $A$; está na
dilatação do complemento quando existe um que o tire --- negação uma da outra, e tomar
complementos troca-as. As peças cuja máscara já é a sua própria antípoda não distinguem os
dois tabuleiros; só o Peão distingue, e move-se \emph{trocando de cor}. Uma promoção de um
lado é uma promoção do outro, lida ao contrário.

E as duas determinam-se uma à outra: $\delta(A)\subseteq Y$ exatamente quando
$A\subseteq\varepsilon(Y)$ --- a mesma frase com os dois quantificadores trocados de
ordem. Daí saem $\delta\varepsilon\delta=\delta$ e
$\varepsilon\delta\varepsilon=\varepsilon$, e daí que repetir não acrescenta:
\textbf{anunciar, confirmar e anunciar de novo é anunciar.} O aperto de mão tem dois
tempos porque um terceiro é absorvido pelo primeiro.

\paragraph{E o mate.} Com $A$ as casas atacadas --- a dilatação de cada atacante pela sua
máscara, com oclusão --- e $O$ as casas dos seus, o rei não tem para onde ir exatamente
quando
\[
k\in\varepsilon(A\cup O).
\]
Dilatar o ataque, erodir com o Rei: \textbf{o predicado da fuga é as duas operações em
sequência.} No mate mais curto do jogo, $1.\,$f3 e5 $2.\,$g4 Dh4, o rei está em e1 e as
cinco casas vizinhas dentro do tabuleiro estão ocupadas pelos seus ou atacadas. O mate
inteiro é mais do que isso --- xeque, sem fuga, não capturar quem dá o xeque, não interpor
---, e as duas últimas condições, nesta formulação, não são morfológicas.

\paragraph{E o que gira contra o que absorve.} O Rei ocupa uma casa: um lance troca duas
coordenadas e deixa o resto, é permutação, e o peso continua um --- nada se junta, nada se
perde. Do outro lado, erodir o tabuleiro inteiro dá $6\times6$, $4\times4$, as quatro
casas centrais, nada; e o meio das quatro é $(7/2,\,7/2)$, que não é casa --- num
tabuleiro de lado ímpar a mesma erosão acaba numa casa só, e a metade vem da paridade, não
do operador. A abertura, essa, apaga: sobre duzentos conjuntos distintos devolve sessenta
imagens. \textbf{O lance reverte; o colapso não} --- é o critério deste capítulo, aqui em
ato. E as duas leituras do tabuleiro, ``só há Reis'' e ``não há Rei'', são a mesma coisa
vista dos dois lados da virada.

\textbf{JOGADA: a torção.} O lance não soma força: \emph{permuta}. Quem joga para juntar
massa perde a reversibilidade --- e o que não tem volta não pertence.

\subsubsection{a régua --- medir sem transbordar}

A unidade é uma régua de comprimento um, e mede-se compondo partes próprias dela: toda
soma parcial fica dentro. Uma soma parcial que passa não é corrigida depois --- com as
parcelas todas positivas, ela já força o total a passar, e a configuração deixa de ser uma
decomposição da unidade. \textbf{É rejeitada na hora.}

\paragraph{E daí sai o ouro.} Contem-se os estados que a régua admite: sequências de
decisões de comprimento $k$, com o padrão $11$ proibido, porque dois passos de saturação
seguidos passariam da unidade. Se a sequência acaba em $0$, os primeiros $k-1$ passos são
qualquer sequência válida; se acaba em $1$, o anterior é forçosamente $0$, e os primeiros
$k-2$ são qualquer sequência válida. Os dois casos não se sobrepõem:
\[
A_k=A_{k-1}+A_{k-2}.
\]
Procurando crescimento da forma $\lambda^k$ e dividindo por $\lambda^{k-2}$, sai
$\lambda^2-\lambda-1=0$ --- a equação do gato com $m=1$. \textbf{O ouro é o que a proibição
do transbordo custa}: livre, a régua daria $2^k$ estados e a taxa seria $\ln2$; com a
proibição, é $\ln\varphi$.

\paragraph{E volta.} A assinatura de $\varphi$ é $[1;1,1,\dots]$, a mais simples que
existe, e as razões que ela regenera são $1,\ 2,\ 3/2,\ 5/3,\ 8/5,\dots$ --- os próprios
$A_k$, um sobre o anterior. A restrição produziu o metal; o metal reproduz a contagem.
\textbf{A régua e o número são o mesmo objeto nos dois sentidos.}

Trocando o passo por $m$, $A_n=m\,A_{n-1}+A_{n-2}$ dá $\lambda^2-m\lambda-1=0$: a família
inteira, chegada pelo outro lado --- contando estados, em vez de compondo matrizes.

\paragraph{E a volta fechada não usa resto.} Sobre uma órbita, a volta completa vale um e
$\theta$ identifica-se com $\theta+1$. Uma soma de incrementos que passa de uma volta
ultrapassa a órbita antes do fechamento, e é rejeitada:

\begin{tcolorbox}[colback=ouroclaro!35,colframe=ouro,boxrule=0pt,leftrule=2pt,arc=0pt,left=4mm,right=3mm,top=2mm,bottom=2mm,breakable]\itshape
$\theta>1$ \textbf{não é reduzido ao resto}. Não é uma posição nova na órbita --- é uma
decomposição inválida. \textbf{Tomar o resto é fingir que o transbordo não aconteceu.}
\end{tcolorbox}

\textbf{JOGADA: a medida.} Um lance que excede a unidade não é aparado nem reciclado ---
ele não aconteceu. Quem corrige o excesso pelo resto joga com um universo que já vazou.

\subsubsection{local e global --- o fechamento como produto}

Um número mede-se no lugar comum --- o valor absoluto --- e em cada lugar primo. Reunidas
todas as medições num objeto só, a soma e o produto correm lugar a lugar.

Esse objeto \textbf{não fecha}. Um elemento inverte quando é não nulo em todos os lugares,
e basta um lugar nulo para não inverter. E há divisor de zero à vista:
\[
(1,0,0,\dots)\cdot(0,1,0,\dots)=0
\]
com ambos os fatores diferentes de zero. Os que são não nulos em toda parte formam um
grupo: \textbf{o anel cinde, e o inverso vive no meio.}

E o metal é o meio do meio. A sua norma é $\sigma_m\sigma_m'=-1$, logo ele inverte, e o
inverso é $\bigl(\sqrt{m^2+4}-m\bigr)/2$. Um inteiro comum não faz isso --- a norma de $2$
é $4$. Cada degrau é mais estreito que o anterior: o anel que cinde, o grupo dos
invertíveis, os que vêm de razões, e as unidades, que são os metais.

\paragraph{E a lei, escrita como produto.} Os que vêm de uma razão satisfazem
\[
\prod_v |x|_v = 1 ,
\]
sobre todos os lugares. \textbf{O que se ganha num, perde-se noutro} --- é a mesma
conservação do reino, agora sem impedância nenhuma, e o fator de fechamento é literalmente
um produto que vale um.

E repare-se no que este corpo mostra, porque é o critério outra vez: ele tem soma,
produto, operador e valuação --- tudo o que se costuma pedir --- e cinde na mesma.
\textbf{Ter as operações não faz o corpo.} Falta a norma ser definida.

\textbf{JOGADA: a alfândega.} Não se mede num lugar só. Um lance é legal quando o produto
de todas as medições dele fecha em um --- e um único lugar zerado sela o resto.

\subsubsection{técnico --- as Princesas, e a régua sem assinatura}

Não produz estado: produz \textbf{veredito}. A régua é o resíduo, e cabe numa linha: $x\oplus x=0$. Ele mede sobre $\mathbb{F}_2$ --- $\oplus$ é a diferença simétrica (o \code{XOR}, o que os dois lados \emph{não} têm em comum), $\otimes$ é a conjunção das hipóteses (o \code{AND}: todas devem valer, não a maioria), $\prod$ é a refutação (o \code{NOT}) ---, e por isso é um dos que não têm assinatura: sobre $\mathbb{F}_2$ ela não existe.

A arquitetura é mínima de propósito. Ali um teorema não é um texto: é um \textbf{sequente} \emph{hipóteses} $\vdash$ \emph{proposição}, de tipo protegido, que só pode ser cunhado por cerca de \textbf{catorze regras primitivas} --- assumir, introduzir e eliminar a implicação, a reflexividade, a simetria, a transitividade. Toda prova do reino, por longa que seja, reduz-se a essas catorze; a confiança concentra-se num núcleo tão pequeno que se lê inteiro numa tarde. E a porta é uma só: passa o par (\emph{testa} $\to$ verdadeiro) $\wedge$ (\emph{refuta} $\to$ falso). Sem refutador, não entra.

E o próprio corpo grava o limite dele: \emph{$x\oplus x=0$ passa nos $k$ casos testados, mas \textbf{a amostra não é o universo}}. O fechamento dele é a \textbf{fase}: o eixo gerar $\bowtie$ atestar \emph{é} o eixo magnitude $\bowtie$ fase --- quem gera dissipa, $|e^{at}|\neq1$; quem atesta conserva, $|e^{i\theta}|=1$. Uma verificação que alterasse o que verifica teria módulo diferente de um, e seria vazamento. \textbf{JOGADA: a atestação} --- resíduo $0$, o lance é legal; resíduo não nulo, o lance \emph{não aconteceu}. A única peça que não ataca, e sem ela nenhum universo abre.

\subsubsection{rotor --- a régua que não satura}

Mede o giro, e a sua régua é hiperbólica: a \textbf{rapidez} $\lambda=\operatorname{artanh} v$, a coordenada do recobrimento. A necessidade dela vem do $\oplus$ deste corpo, que não é a soma comum:
\[
v_1 \oplus v_2 \;=\; \frac{v_1+v_2}{1+v_1v_2}.
\]
Some $0{,}9$ com $0{,}9$: dá $0{,}9945$, não $1{,}8$. A régua tem teto, e o denominador $1+v_1v_2$ é o $\otimes$ escondido dentro do $\oplus$ --- o produto que segura a soma. Então vem Pontryagin na sua forma mais nua e desmancha o teto:
\[
\lambda(v_1\oplus v_2) \;=\; \lambda_1+\lambda_2 .
\]
Não são duas grandezas: é a \textbf{mesma} grandeza em duas réguas. O que satura numa, acumula na outra.

O rotor fecha porque \textbf{conserva}: é a fase, o círculo $|z|=1$, regulador nulo, a órbita que volta ao mesmo ponto. Sozinho é meia coisa; com o motor --- a magnitude, que cresce e dissipa --- fecha um corpo de assinatura $(2,0,0)$: um número é motor vezes rotor, e $\log z=\log|z|+i\theta$ reparte a magnitude (aditiva) da fase (o giro). \textbf{JOGADA: o impulso.} Somar velocidades nunca leva além do limite; somar rapidezes não tem teto. \textbf{A peça que aprende a trocar de régua ganha} --- e é o mesmo lance.

\subsubsection{tabuleiro cósmico --- a escala, e o dipolo}

Não mede grandeza: mede \textbf{expoente}. As densidades somam; o fator de escala compõe por produto, $a(s+t)=a(s)\cdot a(t)$ --- a equação funcional de uma coisa só; e o operador é a expansão
\[
a(t)=e^{Ht},\qquad \log a = H\,t .
\]
Sob o logaritmo, a expansão inteira é uma reta de inclinação $H$. É a mesma tríade do econômico, e não por semelhança: é a mesma conta, porque a expansão \emph{é} o juro composto.

A régua fina é o \textbf{dipolo}: dois atratores, o negro (o poço, o mate) e o branco (o pico, o refúgio), e o potencial de um ponto é a diferença dos inversos das duas distâncias. O flip é a reflexão $\nu=-1$: trocar as fontes leva $V\mapsto-V$, e o atrator vira repulsor. É esse dipolo --- e não uma régua feita para a ocasião --- que mede o gap do rei: o mate de bispo e cavalo é literalmente um problema de dois atratores, com o poço na cor do bispo e o refúgio na oposta. O corpo entrópico é a \textbf{metade} deste: o negro é o máximo de entropia, o branco o mínimo, e o entrópico fica com uma polaridade só. Meia régua.

O fechamento é uma \textbf{lei de potência de resíduo nulo}: endireitada a nuvem, o que sobra é o segundo vetor do endireitamento, e a norma dele é o resíduo --- $3{,}09\cdot10^{-5}$ nos planetas, $6{,}01\cdot10^{-9}$ nas luas. Não é ajuste: é o Princípio da Medida Consistente escrito numa lei de potência. \textbf{JOGADA: a escala.} Um lance cósmico não move a peça: move o \textbf{tamanho do tabuleiro}. Tudo se conserva em proporção --- e quem mede em grandeza se perde, quem mede em expoente não.

\subsubsection{universal --- a raiz, grau 1}

Já construído acima: régua $a^2$, assinatura $(1,0,0)$, operador o sucessor. Resta o que é próprio dele como \emph{instância}, e é um contraste que separa os corpos melhor do que qualquer tabela. Aqui o sucessor \textbf{não} é endomorfismo da soma: $S(2+3)=6$, mas $S(2)+S(3)=7$. No corpo entrópico, onde $\oplus$ é o $\max$, a \emph{mesma} translação \textbf{é} endomorfismo, porque o $\max$ comuta com a translação e fixa o neutro. Mesma operação, dois corpos, dois comportamentos: \textbf{é o corpo que decide, não a fórmula.} \textbf{JOGADA: o passo}, $+1$. Não tem poder nenhum, e toda jogada do reino, decomposta, é uma pilha dele.

\subsubsection{nervoso --- a rede, não o nó}

A novidade é uma negação: os neurônios \textbf{não são atratores isolados}. O que mede não é o nó, é a \textbf{conexão} --- e por isso a matriz deste corpo é a das ligações. A tríade é herdada inteira do eletromagnético: $\oplus$ é a \textbf{integração}, a corrente que chega ao nó e se soma ali, $\sum_i w_ix_i$; $\otimes$ é o \textbf{peso} da ligação, o ganho; e $\prod$ é a \textbf{ativação},
\[
\sigma(x)=\frac{1}{1+e^{-x}},
\]
com o logaritmo da razão do outro lado --- o mesmo $\exp/\log$ dos juros e da expansão, aqui a decidir se o nó dispara.

O corpo nervoso é \textbf{conservativo}, e isso tem marca exata no gerador: os autovalores são \textbf{reais} e os atratores são \textbf{pontos fixos} --- as memórias, os fundos de poço. O seu dual é o motor, que \textbf{dissipa}: ali os autovalores são complexos e os atratores são \textbf{ciclos}. O flip troca real por complexo, e \textbf{a rotação é a marca do dual}. Um guarda, o outro gasta. E a rede de atratores \textbf{é o tabuleiro cósmico em pequena escala}: cada memória um poço local, e a rede desce ao mais próximo. \textbf{JOGADA: a cadeia.} O lance não atinge a peça: atinge \textbf{a rede dela} --- propaga pelos pesos, soma nos nós. Isolar é a defesa; conectar é o poder e o risco.

\subsubsection{exterior --- o produto que se parte em dois}

Aqui o produto revela que sempre foram dois:
\[
ab \;=\; \underbrace{a\cdot b}_{\text{o interior}} \;+\; \underbrace{a\wedge b}_{\text{o exterior}} .
\]
O interior é simétrico e mede o que os dois têm em comum; o exterior é antissimétrico e mede a \textbf{área} que os dois varrem juntos --- e anula-se exatamente quando os lances são paralelos, $a\wedge a=0$. A álgebra é a do elemento nulo, $e^2=0$, com $(a+be)\otimes(c+de)=(ac,\ ad+bc)$.

A régua dele é o que \textbf{falta}: mede pela dimensão que sobra. Em $\R^n$ o $(n+1)$-ésimo resíduo é \textbf{zero} --- a régua tem exatamente $n$ passos, nem um a mais. Ponham-se cinco vetores em $\R^6$ e o resíduo \textbf{não morre}, porque sobra dimensão. Ele não mede o que está lá: mede o buraco. É o único corpo do bestiário com $r>0$ --- assinatura $(1,0,1)$, norma $a^2$ cega à segunda coordenada, $e\neq0$ com $N(e)=0$: divisor de zero, sem inverso, \textbf{não é corpo}. Toda a informação está no \emph{par}, $\langle X(x),X(y)\rangle=-\tfrac12|x-y|^2$, não na norma de cada um --- a mesma régua radical do relógio de xadrez.

E é assim que ele fecha: casado com a fase, o operador que gira ativa a segunda coordenada, $e^2=0$ vira $i^2=-1$, a norma $a^2$ vira $a^2+b^2$ --- e a assinatura sobe de $(1,0,1)$ para $(2,0,0)$. Com $r=0$ e a norma definida, todo elemento não nulo passa a ter inverso. \textbf{Sozinho o exterior é um cone; com o rotor, é um corpo.} \textbf{JOGADA: o corte.} Atacar na direção do inimigo não gera nada; atacar \textbf{transversal} gera a lâmina.

\subsubsection{sensitivo --- a régua hiperbólica, e a escolha}

Assinatura $(1,1,0)$, norma $1-s^2$: uma régua redonda e uma hiperbólica. É literalmente a mesma assinatura do relógio, do expansivo e do evolutivo --- quatro corpos, um instrumento ---, e é a construção que torna isso visível: mesma assinatura, corpos congruentes, a mesma régua vestida de outra base. Os extremos $s=\pm1$ são \textbf{horizontes}: chega-se arbitrariamente perto, nunca se chega; e por isso a composição aqui é a do rotor, $s_1\oplus s_2=(s_1+s_2)/(1+s_1s_2)$.

O casamento deste corpo liga as duas maneiras de medir um mesmo golpe:
\[
\|x\|_\infty \cdot \|x\|_2 = 1 .
\]
Uma régua olha o \textbf{pico}, o maior valor isolado; a outra olha a \textbf{energia}, o todo somado --- e elas são inversas uma da outra. Este é o $\Gamma=0$ do sensitivo, escrito sem impedância nenhuma: não sobra nada entre as duas leituras, porque uma é o recíproco da outra. Daí a lei de quem percebe: \textbf{não dá para ser preciso e forte no mesmo lance.} \textbf{JOGADA: a percepção.} Escolher a régua \emph{é} o lance.

\subsubsection{deflexivo --- o operador que age de fora}

Assinatura $(1,0,0)$ --- a mesma do universal, uma régua só ---, e a régua é a \textbf{reflexão}, $p\mapsto-p$: não se bloqueia, desvia-se, e o que volta volta com o sinal trocado. Mas o que este corpo mostra é outra coisa: \textbf{o operador $A_m$ não pertence ao sistema que o gera. Pertence ao completamento.}

O operador é o gato
\[
A_m=\begin{pmatrix} m & 1\\ 1& 0\end{pmatrix},
\]
de polinômio $x^2-mx-1$, cujas raízes são $\sigma_m$ e $\sigma_m'=-1/\sigma_m$: uma \textbf{expande}, a outra \textbf{contrai}, e o produto é o determinante,
\[
\sigma_m\,\sigma_m'=\det A_m=-1 .
\]
As entradas do gato são inteiras; as raízes, não --- o operador vive fora do sistema de onde saiu. E a deflexão é justamente a troca das duas raízes, $x\mapsto-1/x$, que fixa o produto $-1$ e é o \textbf{flip do esquilo}. É esse par, expandir casado a contrair, o $\Gamma=0$ deste corpo: um boost sem a sua contração seria valor aparecendo fora do Sol.

E aqui o corpo grava o aviso com que este livro fecha, e que vale por todo o bestiário: em régua exata o resíduo vale $1{,}14\cdot10^{-16}$; trocada a fração pelo float, passa a valer exatamente zero. \textbf{O float dá um zero falso.} Com ele vem o segundo aviso, sobre qual órbita se pode habitar: a média temporal do resíduo tende à metade apenas na órbita \textbf{irracional} (na áurea, $0{,}500002$); numa órbita racional como $\theta=1/4$ dá $0{,}375$, que não é metade de nada. \textbf{Só o irracional equidistribui} --- e é por isso que o coração do reino pulsa em $\varphi$, e não numa fração. \textbf{JOGADA: a deflexão.} O operador age \textbf{de fora do tabuleiro} --- o gato.

\subsection{O fecho --- por que o conjunto é um corpo}

Recolha-se a construção inteira, que agora se lê de uma vez.

Há \textbf{um} corpo elementar, o universal: uma régua redonda, grau $1$, e um operador que não foi escolhido --- o sucessor, os caracteres, a órbita da torção. Cada corpo do bestiário sai de uma \textbf{assinatura}: a órbita que colapsa a classe e a regenera, ida e volta, sem terminar. Congelada, ela dá a contagem $(p,q,r)$; a correr, ela dá o gato, e o gato dá o corpo. \textbf{Cada corpo é uma matriz}, e nenhuma opera sozinha --- ligam-se todas por duais, que dão o que falta, e por transformadas, que levam de uma à outra. Reunidas sob $\oplus$ (que soma os graus) e $\otimes$ (La Hire, que compõe), com a torção por invariante e a norma multiplicativa por lei de conservação, essas matrizes formam de novo um único corpo --- o \textbf{Omnitrix} ---, e a torre onde ele vive fecha no grau $8$, assinatura $(8,0,0)$: oito réguas redondas, e não há nona.

E o que o universal faz ali dentro são duas coisas. A \textbf{transformada} leva um corpo a outro, e \textbf{volta}: $\F^{-1}$ é exata, $\F^4=\mathrm{id}$, e ela é unitária --- $\|x\|=\|\F x\|$, nada vaza. Vem em par, porque o corpo tem duas metades: a \textbf{universal} no eixo da fase, a \textbf{dourada} no eixo da taxa, e $\F\circ\log$ leva uma na outra. A \textbf{convolução} funde dois corpos num só, e \textbf{desfaz-se}: o fator que entrou pode ser extraído, exato, sem resíduo. É essa a razão --- a única --- pela qual o conjunto não é uma coleção mas um corpo: \textbf{tudo o que se faz nele, desfaz-se.}

E as inversas não são três sortes separadas. São \textbf{uma} razão, dita três vezes: \textbf{o dual do dual devolve o grupo.} A transformada volta porque ir ao dual e voltar é a identidade; a convolução desfaz-se porque do outro lado ela é um produto, e produto se divide; o inverso de uma matriz é o seu dual sobre a norma, e é ele que decide se aquilo é corpo. Ir e voltar, sempre a mesma frase --- e é dela que sai a proibição de vazar, porque o que não tem volta não é apenas inconveniente: \textbf{não pertence.}

E uma nota sobre a volta, que o livro já apanhou noutro lugar: \textbf{a inversa não é a troca de lados --- é a troca \emph{com o sinal}.} Uma permutação que apenas troca as duas faces \emph{preserva} o telescópio, e um par que preserva onde devia inverter é meia dualidade. A dualidade inteira é a permutação \textbf{com a reflexão}: ela inverte o telescópio, $\sum h\mapsto-\sum h$, e troca a régua pela sua recíproca --- a impedância invertida, que é o mesmo $\exp\leftrightarrow\log$ de sempre. \textbf{Falta-lhe o sinal, e um par sem sinal trocado é o mesmo erro outra vez.} Meia dualidade não fecha; e é justamente sobre o que fecha que estas páginas são.

E é também a lei da alfândega, que este livro enunciou muito antes de chegar aqui: \textbf{passa o que reverte, fica retido o que não tem volta.} Por isso a casa selada arde, e o limite arde: do degrau chega-se ao ponto, do ponto não se recupera o degrau --- a passagem ao limite não tem inversa. Tudo o mais atravessa.

\subsubsection{Identidade atravessa, limite não}

Vale dizer o que essa frase custa na prática, porque é o critério com que tudo neste livro
foi medido.

A cadeia $0{,}9,\ 0{,}99,\ 0{,}999,\dots$ tem termo $1-10^{-n}$: sobe sempre, fica sempre
abaixo de um, e nada abaixo de um a limita. Uma leitura declara que o objeto \emph{é} o
limite, e apaga a subida. A outra guarda a subida --- todo número é a sua volta mais a sua
fase, e a fase vai de zero a um sem chegar; o um não pertence à volta em curso, é o zero
da seguinte. Nessa leitura as truncagens têm todas volta zero, e o um tem volta um:
\textbf{estão em voltas diferentes.} Nenhuma das duas leituras é um erro. Trocá-las uma
pela outra é que é.

E a diferença aparece no metal. Levada aos inteiros, uma \textbf{identidade} continua
identidade: onde a régua de máquina dava $5{,}6\cdot10^{-16}$, a mesma conta em corpo
finito dá $144=144$; a inversibilidade dá bit a bit sem erro; e a linearização deixa de
ser estimar uma curvatura e passa a verificar $\log(c\,t^\alpha)=\log c+\alpha\log t$, que
é igualdade de inteiros no logaritmo discreto --- \emph{não há curvatura a estimar.}

Um \textbf{limite} não atravessa, e há a medição que o mostra. As duas médias --- a
aritmética e a geométrica --- encontram-se na reta porque há \emph{ordem}: a geométrica
fica entre as duas anteriores, e o intervalo aperta. Levadas ao corpo finito, onde a raiz
existe e o encontro poderia ser exato, de setenta e oito pares apenas dois fecham. Não há
ordem, não há entre, não há aperto. Do mesmo objeto, o que atravessa é a parte algébrica
--- o ponto de encontro é invariante, e a média geométrica continua o meio proporcional,
$b'^2=ab$, em milhares de pares sem erro. \textbf{O que não atravessa é a dinâmica.}

\begin{tcolorbox}[colback=ouroclaro!35,colframe=ouro,boxrule=0pt,leftrule=2pt,arc=0pt,left=4mm,right=3mm,top=2mm,bottom=2mm,breakable]\itshape
Antes de medir em resíduo zero, perguntar de que natureza é a coisa. \textbf{Se for um
limite, o resíduo zero ou não existe ou está a ser fabricado.}
\end{tcolorbox}


E fechar não é parar. O que fecha é cada volta --- a transformada que vai e vem, a convolução que funde e se desfaz, o corpo cujo inverso está dentro dele. O que \emph{não} fecha é a subida: há um corpo por assinatura e as assinaturas não acabam; o fractal de tabuleiros toma o nível de baixo como peça e sobe outra vez, e a lei atravessa todos; a assinatura do metal repete o seu passo sem chegar a fim. \textbf{O universo fecha-se em cada volta e segue subindo} --- e é por isso que o mate não mata: a órbita fecha, e o processo continua.

\begin{tcolorbox}[colback=ouroclaro!35,colframe=ouro,boxrule=0pt,leftrule=2pt,arc=0pt,left=4mm,right=3mm,top=2mm,bottom=2mm,breakable]\itshape \textbf{O universal opera no Omnitrix, sobre as matrizes.} Não há vinte e oito mundos: há uma régua, contada com sinal, e duas operações que revertem. \textbf{O gato mede, o esquilo gira, a soma aterra --- e o resíduo é a massa.}\end{tcolorbox}

\vspace{2mm}\hrule\vspace{3mm}

\subsection{O aviso que o corpo deflexivo grava}

De todos os teoremas do catálogo, um deve ser lido antes de qualquer partida:

\begin{tcolorbox}[colback=ouroclaro!35,colframe=ouro,boxrule=0pt,leftrule=2pt,arc=0pt,left=4mm,right=3mm,top=2mm,bottom=2mm,breakable]\itshape \textbf{resíduo exato = \code{1.14e-16}; mutar \code{Fraction} $\to$ \code{float} rende exatamente \code{0.0}}\end{tcolorbox}

O \textbf{float dá um zero falso}. O resíduo verdadeiro não é zero --- o float apenas \textbf{não o enxerga}. É por isso que a
régua do Reino Dourado é o \textbf{exato} ($\mathbb{Q}$, $\mathbb{Z}$[$\varphi$], os inteiros da ISA), e o float só aparece \textbf{no fim}, ao acender o
pixel. Um jogador que confia no zero do float acredita que o seu universo fechou --- e ele está vazando.

E o complemento, também do deflexivo: a média temporal do resíduo tende a \code{?/2} \textbf{apenas na órbita irracional}
(áurea, \code{0.500002}); numa órbita racional (\code{theta = 1/4}) dá \code{0.375 != 1/2}. \textbf{Só o irracional equidistribui} --- é
por isso que o coração do reino pulsa em \code{phi}, e não numa fração.

\subsection{E tudo isto cabe em dez bytes}

A construção acabou, e ela correu na reta: uma raiz, uma assinatura, matrizes,
duas operações que revertem.

Nada disso precisa da reta para acontecer. Precisa de dois inteiros, cinco
opcodes e uma memória --- e é o que a parte seguinte escreve, byte a byte, de
modo que se possa reconstruir sem voltar a estas páginas.

\textbf{Não é uma aplicação do que se derivou.} É a mesma coisa outra vez, e a
única diferença é que ali ela tem endereço. O gato, que aqui gerou o corpo, é lá
o que a memória faz ao ler; as duas projeções dele, que aqui deram $\oplus$ e
$\otimes$, são lá o \emph{byte da operação}; a torção, que aqui deu os
caracteres, é lá um número com cinco dígitos; e a transformada que aqui volta
sem perder volta lá pela mesma razão e com a mesma normalização, escrita em
inteiros.

\begin{tcolorbox}[colback=ouroclaro!35,colframe=ouro,boxrule=0pt,leftrule=2pt,arc=0pt,left=4mm,right=3mm,top=2mm,bottom=2mm,breakable]\itshape
Quem leu a construção já leu a máquina. \textbf{Falta ver os bytes.}
\end{tcolorbox}


\subsection{A zeta dinâmica e os problemas em aberto --- onde a analogia para}\label{sec:clay}

A zeta deste texto, $\zeta(x) = 1/\beta^*(x)$, é do mesmo \emph{tipo} que as zetas para as quais a
Hipótese de Riemann está \textbf{provada} --- as de variedades sobre corpos finitos, conjecturadas
por Weil e demonstradas por Deligne em 1974. Ambas são \emph{racionais}, ambas têm equação
funcional. Vale a pena percorrer a analogia até ela parar, porque o sítio onde para é informativo.

\paragraph{A equação funcional, essa existe.} O reverso de $\beta^*$ devolve $-\beta$: para
$(n,m)=(2,1)$, $x^2\beta^*(1/x) = -x^2+x+1 = -\beta$; para $(3,2)$, $x^3\beta^*(1/x) =
-x^3+2x^2+1$. \emph{É a involução $\nu$ a relacionar os dois lados} --- e é isso, precisamente, uma
equação funcional.

\paragraph{Mas a Hipótese de Riemann não tem análogo aqui.} Nas zetas de Weil, o enunciado é que
os zeros têm \textbf{todos} o mesmo módulo, $q^{-1/2}$ --- é essa a ``recta crítica'' do caso
finito, e é ela que Deligne provou. Na nossa, os polos têm módulos \emph{diferentes}:

\begin{center}\small
\begin{tabular}{ll}
\hline
$(n,m)$ & módulos dos polos \\ \hline
$(2,1)$ & $0{,}61803$; $1{,}61803$ \\
$(3,2)$ & $0{,}45340$; $1{,}48512$; $1{,}48512$ \\
$(4,3)$ & $0{,}32941$; $1{,}40419$; $1{,}40419$; $1{,}53961$ \\
\hline
\end{tabular}
\end{center}

\noindent Seria fácil arquivar isto como ``não há analogia''. Seria o erro de medir um lado só:
\textbf{pureza e dominância não são presença e ausência --- são os dois extremos do mesmo eixo}, e
esta zeta está num deles.

\begin{center}\small
\begin{tabular}{lll}
\hline
& os módulos dos zeros & o que daí resulta \\ \hline
Weil (provado por Deligne) & todos iguais: \emph{puros} & equidistribuição \\
Pisot (Teor.~o teorema de Pisot para $\beta_{n,m}$ no \emph{Catálogo}) & um domina, os outros contraem & \textbf{concentração} \\
\hline
\end{tabular}
\end{center}

\noindent E o lado que falta medir não é o aditivo --- é o \emph{multiplicativo}. A propriedade de
Pisot não é sobre $\{k\sigma\}$ (a órbita dos múltiplos, que a a equidistribuição em $\R^{n}$ no \emph{Catálogo} mede) mas sobre
$\|\sigma^k\|$, a distância de cada \emph{potência} ao inteiro mais próximo. Medido, e os dois lados
são opostos:

\begin{center}\small
\begin{tabular}{rrr}
\hline
$k$ & $\|\varphi^k\|$ (potências) & $\|k\varphi\|$ (múltiplos) \\ \hline
$5$ & $9{,}0\times10^{-2}$ & $0{,}0902$ \\
$10$ & $8{,}1\times10^{-3}$ & $0{,}1803$ \\
$20$ & $6{,}6\times10^{-5}$ & $0{,}3607$ \\
$40$ & $3{,}0\times10^{-7}$ & $0{,}2786$ \\
\hline
\end{tabular}
\end{center}

\noindent A órbita das potências \textbf{colapsa} no inteiro; a dos múltiplos espalha-se. E o teste
distingue: a média de $\|x^k\|$ para $k\leq40$ dá $0{,}035$ em $\varphi$ e $0{,}050$ em $1+\sqrt2$
(ambos de Pisot), contra $0{,}125$ em $\sqrt2$ e $0{,}177$ em $\pi$, que não o são. (Para $k$ grande
os valores voltam a subir --- isso é o \emph{epsilon} do cálculo em vírgula flutuante, não o
fenómeno.)

\emph{E a não-pureza é o que dá a ordem.} É porque $|\lambda_2|/\sigma < 1$ --- $0{,}382$ em
$(2,1)$, $0{,}305$ em $(3,2)$ --- que a série converge, que os traços são inteiros
(Prop.~os traços no \emph{Catálogo}) e que a base ordena. \textbf{Se os módulos fossem puros, nada disso
existiria.} O que aqui não há é a \emph{pergunta} de Weil; o que há é a resposta oposta, e ela é
que sustenta o resto do texto.

\begin{obs}[e os sete problemas, sem inflacionar]\label{obs:clay}
Convém arrumar isto de uma vez, porque a tentação de reclamar proximidade é grande:

\begin{center}\small
\begin{tabular}{lll}
\hline
problema & toca este trabalho? & como \\ \hline
Hipótese de Riemann & só na linguagem & a recta crítica é conjunto fixo (Obs.~a crítica ao critério no \emph{Catálogo}) \\
Birch--Swinnerton-Dyer & tangencialmente & curvas elípticas, via a discussão de Wiles \\
$P$ vs $NP$ & não & --- \\
Navier--Stokes & não & --- \\
Yang--Mills & não & --- \\
Hodge & não & --- \\
\hline
\end{tabular}
\end{center}

\noindent \textbf{Um de sete toca, e toca só na linguagem.} O que este texto tem sobre a recta
crítica é que ela é o conjunto fixo de $s\mapsto1-\bar s$ --- facto sobre a \emph{aplicação}, que
seria verdade ainda que $\zeta$ não existisse. Não é um passo em direcção à Hipótese; é a
observação de que o padrão da casa, ali, não morde. Dizer mais do que isso seria inventar.
\end{obs}


\section{A definição de $\pi$, e a sua construção}\label{sec:pi}

\noindent $\pi$ põe a arquitetura deste texto à prova, e vale a pena dizer porquê antes de o
definir. Por Lindemann, $\pi$ não é raiz de nenhum polinómio de coeficientes inteiros: \textbf{não
está em nenhum dos anéis que a Parte~I constrói}. Se a sua definição tivesse de ser geométrica ---
comprimento de arco, área de círculo ---, então o corpo dual precisaria da geometria para se
completar, e a ordem das partes deste texto estaria trocada.

\medskip
\noindent Não precisa. A definição é \emph{dinâmica}, e sai das duas leis sem pedir nada de novo:
a Lei~$2$ fornece o gerador, a exponencial fornece o fluxo, e a Lei~$1$ fornece o alvo.

\begin{mdframed}[linewidth=0.4pt,innerleftmargin=10pt,innerrightmargin=10pt,
                 innertopmargin=8pt,innerbottommargin=8pt,skipabove=6pt,skipbelow=6pt]
\textbf{Definição.} Seja $J$ um gerador anti-autoadjunto, $J^{\dagger}=-J$, e $\exp(tJ)$ o fluxo
que ele gera. Então
\[
  \boxed{\;\pi \;:=\; \min\{\,t>0 \;:\; \exp(tJ)\cdot 1 = -1\,\}\;}
\]
isto é: \textbf{$\pi$ é o tempo que o fluxo leva a realizar a Lei~$1$.}
\end{mdframed}

\noindent Nenhuma palavra desta frase é geométrica. Não há círculo, não há arco, não há área ---
e o círculo, se aparecer, aparece \emph{depois}, como a figura que este fluxo desenha. É a mesma
inversão que atravessa o texto todo: \emph{a figura não é o dado, é o rasto}.

\subsection{Por que $\pi$ não é um elemento, e sim um limite}\label{sec:pilimite}

\noindent A álgebra dá os passos e não dá o alvo, e isto pode ser medido em vez de afirmado. Os
passos que a álgebra alcança são as rotações de coordenadas racionais --- as ternas pitagóricas
---, e nelas a Lei~$2$ vale \emph{exatamente}: a segunda coordenada da cruz é $1$, sem
aproximação. Mas o alvo da Lei~$1$, que é traço $-2$, nunca é atingido por um passo próprio.

\medido{$2\,376$ pontos racionais do círculo varridos: os $400$ que realizam traço $-2$ são
\emph{todos} degenerados (a segunda componente é nula, e isso é a paragem, não um passo). Nenhum
passo próprio leva a unidade ao seu dual. E a órbita aproxima-se sem chegar: o traço mais próximo
alcançado é $-3968/1985$, estritamente maior que $-2$.}

\noindent \textbf{Daí a construção.} Um objeto que a órbita aproxima sem atingir constrói-se como
o \emph{par} que o encaixota --- que é a cruz outra vez, agora com uma coordenada por baixo e
outra por cima. \emph{$\pi$ não é um elemento do corpo: é a cruz de uma órbita que não fecha.}

\medido{\code{tests/pidual.c} --- $12$ unidades, resíduo $0$, e \textbf{sem vírgula flutuante}:
$\pi$ é produzido por três somas alternadas independentes em aritmética inteira (Machin, Euler,
Hermann), que concordam nas $24$ primeiras casas. Nenhum dígito foi escrito à mão --- se um dos
caminhos estivesse errado, a asserção cairia.}

\subsection{E o alvo da Lei 1 é um ponto estacionário}\label{sec:piestacionario}

\noindent Este parágrafo existe porque a medição o produziu, e não estava previsto. Ao tentar
fazer a verificação \emph{falhar} --- perturbando $\pi$ para ver o resíduo subir ---, o resíduo
não se moveu. Não é defeito da medida: \textbf{$\pi$ é ponto estacionário de $\cos$}, a derivada
lá é $-\sin\pi = 0$, e portanto $\cos\pi = -1$ é insensível a erro de primeira ordem. Verificar a
definição por essa coordenada é um teste fraco.

\medskip
\noindent E o facto tem conteúdo próprio, que é a razão de ficar escrito: \emph{o alvo da Lei~$1$
é um ponto estacionário do fluxo}. A meia volta é estável --- chega-se lá e não se derrapa. É o
mesmo que dizer, noutras palavras, que a órbita não dissipa; e é por isso que a involução é uma
operação e não um acidente.

\medido{$\cos\pi = -1$ com resíduo \emph{zero} unidades contra um limite de $104$ derivado do
próprio algoritmo (divisões inteiras contadas, mais o critério da série alternada --- nenhum
limiar escolhido à mão). Já $\sin\pi = 0$ dá resíduo $3$, e perturbar $\pi$ na décima casa
leva-o a $100\,000$: \textbf{a segunda coordenada separa por cinco ordens de grandeza onde a
primeira não separa de todo.} Mais uma vez, uma coordenada só não chega.}

\subsection{E a razão entre os dois períodos}\label{sec:piperiodo}

\noindent Fica por dizer onde $\pi$ se encaixa no resto. A estaca tem período $2$; o fluxo tem
período $2\pi$. \textbf{$\pi$ é o que traduz um período no outro} --- é a meia volta com
parâmetro, onde a estaca é a meia volta de um só golpe.

\medskip
\noindent Isto arruma uma observação que de outro modo pareceria folclore: $\pi$ aparece em toda a
parte onde há uma involução a ser feita \emph{continuamente}, e não aparece em sítio nenhum onde a
involução seja feita de uma vez. Não é uma constante da geometria que invadiu a análise --- é
\textbf{o preço de fazer em tempo contínuo aquilo que a estaca faz num passo}.


\section{E o gancho: o operador geométrico é o tempo}\label{sec:gancho}
\hfill \emph{\small[navegação: fecha a Análise e abre a Geometria]}

Chegados aqui, a Análise deu três coisas e elas encaixam numa só. Deu a \textbf{métrica} --- e com
ela o limite, a continuidade e a derivada. Deu o \textbf{corpo cruz} e o \textbf{corpo de corpos}
--- onde a medida passa a ser sobre o conjunto e não sobre um elemento. E deu a \textbf{dinâmica}:
a Lei~$2$ em forma diferencial, $(d/dt)^{\dagger}=-d/dt$, que é o que torna o tempo reversível.

\paragraph{E é aqui que a ordem e o tempo se separam, e a distinção abre a Geometria.} A ordem da
torre é o \textbf{relógio}: estática, marca antes e depois, e sobe de andar em andar pela involução.
O \textbf{tempo} é a \emph{dinâmica} desse relógio --- o que ele faz ---, e quem a fixa é a
\textbf{assinatura da régua}: a matriz que a gera. Duas réguas com assinaturas diferentes marcam a
mesma ordem e correm tempos diferentes.

\medskip
\noindent \textbf{E daí o gancho, que é o enunciado do capítulo seguinte:}

\begin{center}\small
\begin{minipage}{0.9\textwidth}
\emph{O operador geométrico é o tempo.} A curvatura não é uma propriedade acrescentada ao espaço:
é o que o gerador da dinâmica faz ao ser aplicado duas vezes --- $f''=K f$ ---, e o seu sinal
decide a geometria. \textbf{Medir o espaço é ver o que o tempo lhe fez.}
\end{minipage}
\end{center}

\noindent É essa identificação que a \emph{Geometria} desenvolve: a equação de Jacobi como a mesma
lei, $\Delta=\operatorname{tr}^{2}-4\det$ como a curvatura de cada corpo, e o expoente da gravitação
a sair da dimensão.

\part{Geometria --- a escrita vista}
\noindent\hfill\emph{\small[\textbf{realização}: a Lei~2 em figura, e o operador é o tempo]}\medskip

\section{A geometria sai das duas leis --- e o texto já a tinha, dispersa}\label{sec:geo}
\hfill \emph{\small[consequência: deriva-se das duas leis, e não abre lugar novo]}

As três partes anteriores repartem-se pelas duas leis, e não arbitrariamente: \emph{a Análise é a
\textbf{leitura}} --- mede, conta, distribui, e a Lei~$1$ é a sua ---, e \emph{a Álgebra é a
\textbf{escrita}} --- gera, opera, constrói, e é da Lei~$2$. A Topologia sobe de uma e a Álgebra
desce da outra. Falta a quarta, e ela é a escrita \emph{vista}: \textbf{a Geometria}.

\noindent Nada aqui é material novo. As peças estão nas três partes anteriores --- a curvatura, o
cone, Euler, a projetiva, a codimensão, a realização hiperbólica --- e o que esta parte faz é
\emph{derivá-las das duas leis} e pô-las lado a lado.

\subsection{A curvatura é o discriminante da borda}\label{sec:geocurv}

Escrita com uma constante no lugar do sinal, a Lei~$2$ é $f''=K\,f$, e $K$ \textbf{é a curvatura}.
Da característica $x^{2}-\operatorname{tr}x+\det=0$ vem
\[
  \Delta = \operatorname{tr}^{2}-4\det ,
\]
e o sinal de $\Delta$ diz a geometria --- que é, palavra por palavra, a classificação clássica das
transformações de Möbius:

\begin{center}\small
\begin{tabular}{@{}ccll@{}}
\toprule
$\Delta$ & a equação & o invariante & a geometria \\
\midrule
$<0$ & $f''=-f$ & $\cos^{2}+\sin^{2}=1$   & elíptica --- o rotor, o Pégaso \\
$=0$ & $f''=0$  & a reta                  & parabólica --- e é aqui que vive o potencial \\
$>0$ & $f''=+f$ & $\cosh^{2}-\sinh^{2}=1$ & hiperbólica --- onde vive $\sigma$ \\
\bottomrule
\end{tabular}
\end{center}

\noindent \emph{Os três invariantes têm o mesmo $1$ do lado direito}; o que os separa é só o sinal,
que é $\sigma\sigma'=\mp1$ --- a alfândega de \S\ref{sec:enunciado}, agora em geometria. As três
geometrias de curvatura constante são as únicas simplesmente conexas em cada
dimensão~\cite{do-carmo}, e é por isso que a lista tem três linhas e não mais. \emph{É o
Programa de Erlangen lido ao contrário}~\cite{klein1872}: em vez de classificar geometrias pelo grupo
que as preserva, lê-se o grupo na borda e a geometria sai dele. E ``curvatura constante'' e ``não
dissipa'' são a mesma exigência: \emph{o que não varia ao longo da órbita é o que não varia de
ponto para ponto}.

\subsection{E daí a curvatura de cada corpo}\label{sec:geocorpos}

Sendo $\Delta$ lido na assinatura, cada corpo deste texto tem a sua, sem que seja preciso
acrescentar-lhe nada:

\begin{center}\small
\setlength\tabcolsep{5pt}
\begin{tabular}{@{}lrrrl@{}}
\toprule
a peça & $\operatorname{tr}$ & $\det$ & $\Delta$ & a curvatura \\
\midrule
$J$ --- o espelho, $J^{2}=+I$ & $0$ & $-1$ & $4$ & hiperbólica \\
$i$ --- a rotação, $i^{2}=-I$ & $0$ & $+1$ & $-4$ & \textbf{elíptica} \\
$M$ --- o gerador de Fibonacci & $1$ & $-1$ & $5$ & hiperbólica \\
o \emph{shift} & $2$ & $+1$ & $0$ & \textbf{parabólica} \\
a família metálica & $m$ & $-1$ & $m^{2}+4$ & hiperbólica, \emph{sempre} \\
\bottomrule
\end{tabular}
\end{center}

\noindent \textbf{O espelho vive na hipérbole e a rotação no círculo}, e agora há um número que o
diz --- e é o mesmo número que decide a ordem em os corpos $K_{n,m}$ no \emph{Catálogo} e a irracionalidade em
onde Pisot falha no \emph{Catálogo}. A classificação por $\Delta$ é clássica~\cite{coxeter,berger}; o que aqui se
acrescenta é que ela se lê \emph{na assinatura do corpo}, sem sair da álgebra. \emph{E uma distinção que a medição impôs}: em $m=0$ o discriminante é $4=2^{2}$, quadrado
perfeito --- e está certo, porque é o nível $0$ e as raízes $\pm1$ são racionais. Hiperbólico vale
para todo $m$; irracional só de $m=1$ em diante.

\subsection{A separação é a codimensão --- e daí a gravitação}\label{sec:geograv}

Do lado da leitura, a Lei~$1$ dá o que separa. Num ambiente de dimensão $n$, remover um subconjunto
de dimensão $k$ separa-o \emph{se e só se} $n-k=1$ (no \emph{Catálogo}): \textbf{codimensão $1$
separa, codimensão $\geq2$ não}. É o índice que a dualidade de Poincaré produz, e é ele que impede a
bijeção contínua $[0,1]\to[0,1]^{n}$.

\noindent E é o mesmo índice que dá a força. Fora da fonte, a Lei~$2$ na forma nula é
$\nabla^{2}\varphi=0$; em simetria radial, $r^{\,d-1}\varphi'=\text{const}$, donde
\[
  F \propto \frac{1}{r^{\,d-1}}, \qquad\text{e em } d=3: \quad F\propto\frac{1}{r^{2}} .
\]
\textbf{O expoente é a codimensão da esfera que envolve a fonte}, e o fluxo $F\times\text{área}$ é
constante porque nada se perde entre esferas. \emph{A órbita que não dissipa e o inverso do quadrado
são a mesma afirmação em variáveis diferentes}: a norma não se move no \emph{tempo}, o fluxo não se
move no \emph{raio}.

\subsection{As faces geométricas da dualidade, e o que fecha cada uma}\label{sec:geofaces}

As mesmas duas leis dão as formas geométricas do par, e em cada uma há um invariante que a troca não
move --- \emph{é ele que faz a volta ser possível}:

\begin{center}\small
\begin{tabular}{@{}llll@{}}
\toprule
a face & o dual & o bidual & o que a fecha \\
\midrule
o cone      & $K^{*}$   & $K^{**}=K$ & a convexidade (o bipolar) \\
o poliedro  & $V\leftrightarrow F$ & o próprio & a aresta, que $\chi=2$ conserva \\
o plano projetivo & ponto $\leftrightarrow$ recta & o ponto & a incidência \\
a variedade & $H_{n-k}$ & $H^{k}$ & a classe fundamental \\
\bottomrule
\end{tabular}
\end{center}

\noindent A bidualidade do cone é o teorema do bipolar~\cite{rockafellar}, e a da variedade é
Poincaré~\cite{thurston}. E o cone dá o aviso que vale para todas: \textbf{ele só é autodual na
métrica dele}. Com o
produto euclidiano emprestado, o dual sai maior --- $111/91/71$ contra $203/223/245$ --- que é o
\emph{preço da régua errada}, agora em geometria.

\medido{Curvatura, codimensão e as quatro faces em \code{tests/bidual.c} ($28$ unidades, resíduo
$0$); o cone e a projetiva em \code{tests/pontofixo.c}; a realização hiperbólica com
$\ell=4\log\sigma$ na Parte II.}

\subsection{Três pontes clássicas, e o que cada uma liga}\label{sec:geopontes}

A geometria riemanniana tem três resultados que ligam, cada um, dois lados que este texto já tinha
\emph{separados}. Não são resultados novos --- são clássicos, e vão citados; o que aqui se faz é
mostrar em que ponto do texto cada um entra.

\paragraph{Primeira: a equação de Jacobi \emph{é} a Lei 2.} O desvio geodésico --- quanto duas
geodésicas vizinhas se afastam --- é governado por
\[
  J'' + K\,J = 0 ,
\]
isto é $J''=-K\,J$~\cite{do-carmo,cheeger-ebin}. \textbf{É a mesma equação} da
\S\ref{sec:geocurv}, com $J$ no lugar de $f$. O que ali se chama \emph{campo de Jacobi} é aqui a
solução da lei, e o que ali se chama \emph{curvatura seccional} é aqui o $\Delta$ da borda. E a
trichotomia sai por outro caminho e dá o mesmo:

\begin{center}\small
\begin{tabular}{@{}cccl@{}}
\toprule
$K$ & $J(t)$ & as geodésicas & a geometria \\
\midrule
$>0$ & $\sin t$  & \textbf{convergem} --- reencontram-se & esférica \\
$=0$ & $t$       & ficam paralelas                        & plana \\
$<0$ & $\sinh t$ & \textbf{divergem}, exponencialmente    & hiperbólica \\
\bottomrule
\end{tabular}
\end{center}

\paragraph{Segunda: Gauss--Bonnet é uma ponte \emph{entre as duas leis}.} Numa superfície fechada,
\[
  \iint_{M} K\,dA \;=\; 2\pi\,\chi(M) .
\]
À esquerda está a \textbf{curvatura} --- geometria, a escrita vista, Lei~$2$. À direita está
$\chi=V-E+F$ --- \textbf{contagem}, a leitura, Lei~$1$ (Euler e o miolo no \emph{Catálogo}). \emph{Este texto tinha
os dois lados e não os tinha ligado}: a curvatura entrou pela borda, e Euler entrou pela
dualidade dos poliedros, por caminhos que nunca se cruzaram. \textbf{Gauss--Bonnet é o cruzamento},
e é o que garante que a nossa Lei~$2$ em figura e a nossa Lei~$1$ em contagem falam do mesmo
objeto.

\medido{Na esfera de raio $R$: $K=1/R^{2}$ e a área é $4\pi R^{2}$, logo $\iint K\,dA=4\pi$
\textbf{para todo $R$} --- derivado, e igual a $2\pi\chi$ com $\chi=2$. E $\chi=2$ é o mesmo que os
cinco sólidos platónicos dão, medido em \code{tests/bidual.c}.}

\paragraph{Terceira: a curvatura decide se o espaço fecha --- e aqui a leitura é nossa.} Bonnet e
Myers: curvatura $\geq\delta>0$ obriga a variedade a ser \textbf{compacta}, com diâmetro
$\leq\pi/\sqrt\delta$~\cite{myers1941}; e Cartan--Hadamard: curvatura $\leq0$ dá recobrimento
universal $\R^{n}$, \emph{aberto}~\cite{do-carmo}.

\smallskip
\noindent \emph{E é aqui que é preciso dizer o estatuto}, porque o que se segue é uma analogia
nossa e não um teorema. O texto mediu que a bidualidade \emph{fecha} quando o invariante cobre o
objeto inteiro, e que em dimensão infinita o bidual é \emph{estritamente maior}
(\S\ref{sec:bidualfaces}). A dicotomia geométrica tem a mesma forma:

\begin{center}\small
\begin{tabular}{@{}lll@{}}
\toprule
curvatura & o espaço & e no texto \\
\midrule
$K>0$ & compacto --- \emph{fecha} & o invariante cobre o objeto inteiro \\
$K=0$ & plano, aberto & a fronteira --- e é onde vive o potencial \\
$K<0$ & aberto, $\R^{n}$ & o bidual é maior: a memória não volta inteira \\
\bottomrule
\end{tabular}
\end{center}

\noindent \textbf{Mas é analogia de forma, e não identidade}: Bonnet--Myers fala de variedades e a
nossa medição fala de espaços vetoriais e anéis, e não há aqui teorema que passe de um para o
outro. \emph{O que se afirma é que a mesma dicotomia --- fechar ou abrir --- aparece nos dois, e
que nos dois o que decide é se há uma quantidade positiva que não se anula.}

\subsection{A régua curva-se: métrica, transporte e geodésica}\label{sec:geometrica}

Até aqui a geometria entrou pelo \emph{sinal} --- que curvatura, que classe. Falta o que ela dá de
ferramenta, e é disso que as aplicações precisam: \textbf{como se mede, como se move, e qual é o
caminho que não gasta}.

\paragraph{A métrica é a régua do próprio objeto, e não uma escolha.} Num espaço curvo o produto
interno muda de ponto para ponto: é o \emph{tensor métrico} $g$. E onde o objeto é uma
\emph{distribuição}, $g$ não se escolhe --- deduz-se. Para a Bernoulli de parâmetro $p$,
\[
  g(p) \;=\; \frac{1}{p\,(1-p)} ,
\]
que \textbf{explode nas pontas e é mínima em $p=1/2$}. Distinguir $0{,}5$ de $0{,}51$ é barato;
distinguir $0{,}99$ de $0{,}999$ é caro. \emph{É a mesma diferença euclidiana e não é a mesma
distância} --- e é isso, e só isso, que a palavra ``curvatura'' quer dizer quando o objeto é uma
distribuição.

\medido{$g(p)=100/\bigl(k(10-k)\bigr)$ em racionais exatos para $p=k/10$: o mínimo é $4$, em
$p=1/2$, \emph{derivado e não escrito} (\code{tests/bidual.c}). E a distância geodésica de Fisher,
$d(p,q)=2\,\lvert\arcsin\sqrt p-\arcsin\sqrt q\rvert$, dá $2{,}0\times$ a euclidiana perto de $1/2$
e $15{,}2\times$ entre $0{,}99$ e $0{,}999$.}

\paragraph{O transporte paralelo é mover sem perder --- e no círculo é o nosso $i$.} Para comparar
dois vetores em pontos diferentes é preciso levar um até ao outro sem o deformar: é o
\emph{transporte paralelo}, e a sua condição é conservar o produto interno. No círculo, transportar
por um quarto de volta é
$\bigl(\begin{smallmatrix}0&-1\\1&0\end{smallmatrix}\bigr)$ --- \textbf{exatamente o $i$ da
fatorização $M=J\cdot i$}. \emph{A rotação que ordena e o transporte que não perde são a mesma
matriz.}

\medido{Norma preservada em $289$ vetores de entradas inteiras, resíduo $0$.}

\paragraph{E a geodésica é o caminho que não gasta.} É a curva cujo vetor tangente é transportado
paralelamente a si próprio --- \emph{anda sem virar}. Nas três curvaturas ela é o que o
\S\ref{sec:geopontes} já deu pelas soluções de Jacobi; e na realização hiperbólica deste texto o seu
comprimento é $\ell=4\log\sigma$~\cite{series1985}: \textbf{o logaritmo da borda é uma distância}.

\subsection{As conexões duais: a nossa estrutura, com o nome dela}\label{sec:geoamari}

E há um sítio onde tudo isto já está escrito com a nossa gramática, sem nos conhecer. Numa
variedade com métrica, chama-se \emph{conexão dual} de $\nabla$ à única $\nabla^{*}$ que
satisfaz~\cite{amari-nagaoka}
\[
  X\,g(Y,Z) \;=\; g(\nabla_{X}Y,\,Z) \;+\; g(Y,\,\nabla^{*}_{X}Z) .
\]
\textbf{É a adjunção}, palavra por palavra: comparar com
$\langle Ax,y\rangle=\langle x,A^{\top}y\rangle$ e ver que a única diferença é o nome. \emph{Mover a
operação de um lado do produto para o outro troca-a pela sua dual.}

\paragraph{E a família tem ponto fixo, exatamente onde o nosso enunciado o põe.} As
$\alpha$-conexões formam uma família em que $\nabla^{(\alpha)}$ e $\nabla^{(-\alpha)}$ são duais uma
da outra, e
\[
  \nabla^{(0)} = \text{a conexão de Levi-Civita} \quad\text{é \textbf{autodual}: } \nabla=\nabla^{*} .
\]
A involução é $\alpha\mapsto-\alpha$; a bidualidade devolve o próprio; e o \emph{único} ponto fixo é
$\alpha=0$. \textbf{É o desenho da nossa fronteira, ponto por ponto}: um par trocado por involução,
e um sítio no meio onde os dois lados deixam de se distinguir (\S\ref{sec:enunciado}).

\medido{$\alpha\mapsto-\alpha$ em $41$ valores: bidual devolve o próprio sem falha, e o ponto fixo
é \textbf{um só} (\code{tests/bidual.c}).}

\begin{obs}[o que isto é e o que não é]
\emph{Não se afirma que a geometria da informação seja um caso deste texto}, nem o contrário. O que
está medido é que a \textbf{mesma estrutura} --- par dual, involução, ponto fixo autodual --- aparece
nos dois, e que ali ela já tinha nome próprio: \emph{conexão dual}, e até
\emph{bidualidade}~\cite{amari-nagaoka}. \textbf{É uma confirmação de que o desenho não é
particular deste objeto}, e é para isso que serve.
\end{obs}

\subsection{\texorpdfstring{$\exp$ e $\log$}{exp e log}: o par aditivo/multiplicativo com ponto-base}\label{sec:geoexplog}

O par que este texto usa desde o princípio --- somar de um lado, multiplicar do outro --- tem em
geometria uma forma com \emph{ponto-base}, e é a que as aplicações usam:
\[
  \exp_{p}: T_{p}M \to M \quad\text{(a soma vira trajetória)},
  \qquad
  \log_{p}: M \to T_{p}M \quad\text{(a trajetória vira soma)},
\]
inversas uma da outra, $\log_{p}(\exp_{p}v)=v$. \emph{O espaço tangente é plano e é onde se soma; a
variedade é curva e é onde o objeto vive.} No espaço plano $\exp$ é a identidade --- somar \emph{é}
andar ---, e \textbf{a diferença entre os dois é exatamente a curvatura}.

\paragraph{E nas matrizes definidas positivas, que é onde vivem os encaixes.} Ali
$\exp_{I}(V)=e^{V}$ e $\log_{I}(P)=\ln P$, e a distância geodésica é
$\lVert\ln(P^{-1/2}QP^{-1/2})\rVert$. A consequência é prática e mede-se:

\begin{center}\small
\begin{tabular}{@{}rrrrr@{}}
\toprule
$a$ & $b$ & euclidiana & geodésica & razão \\
\midrule
$1$    & $2$  & $1{,}00$ & $0{,}69$ & $0{,}69\times$ \\
$1$    & $10$ & $9{,}00$ & $2{,}30$ & $0{,}26\times$ \\
$0{,}1$ & $1$  & $0{,}90$ & $2{,}30$ & $\mathbf{2{,}56\times}$ \\
$0{,}01$& $1$  & $0{,}99$ & $4{,}61$ & $\mathbf{4{,}65\times}$ \\
\bottomrule
\end{tabular}
\end{center}

\noindent \emph{Na escala pequena a geodésica é maior; na grande é menor.} Ir de $0{,}1$ a $1$ custa
$2{,}30$ e não $0{,}90$ --- e ir de $1$ a $10$ custa o mesmo $2{,}30$, embora a diferença euclidiana
seja dez vezes maior. \textbf{A régua multiplicativa não é a aditiva}, e é por isso que somar
encaixes e multiplicar razões dão respostas diferentes: são operações em \emph{lados diferentes} do
mesmo par.

\paragraph{E é isto que as aplicações precisam.} Quando o que se compara são distribuições --- ou
representações que se comportam como tais --- \emph{a distância euclidiana é a régua errada}, e o
erro cresce onde mais importa: nas pontas, onde a probabilidade satura. A métrica de Fisher dá a
régua do objeto; o transporte paralelo dá como comparar dois pontos distantes; e a geodésica dá o
caminho de custo mínimo. \emph{São as três coisas que um encaixe precisa de ter para não medir a sua
própria convenção} (o corpo métrico no \emph{Catálogo}).

\subsection{As quatro partes são autoduais duas a duas}\label{sec:geoarvore}

As quatro não são uma lista: são \emph{dois pares}, e o conjunto é fechado sob a troca.

\begin{center}\small
\begin{tabular}{@{}lll@{}}
\toprule
 & \textbf{escrita} --- Lei~$2$ & \textbf{leitura} --- Lei~$1$ \\
\midrule
\emph{opera} & \textbf{Álgebra} --- gera, $+$, $\times$, inverso & \textbf{Análise} --- mede, taxa, distribuição \\
\emph{vê}    & \textbf{Geometria} --- a mesma lei em figura      & \textbf{Topologia} --- ordem, limite, vizinhança \\
\bottomrule
\end{tabular}
\end{center}

\noindent \emph{Dualizar troca as colunas e devolve o quadro} --- é o que quer dizer serem
autoduais duas a duas: Álgebra $\leftrightarrow$ Análise (o que opera), Geometria
$\leftrightarrow$ Topologia (o que vê).

\paragraph{E a árvore de dependências não foi desenhada: foi medida.} Contando as referências
cruzadas entre partes, $70$ ao todo:

\begin{center}\small
\begin{tabular}{@{}lrrrr@{}}
\toprule
de $\backslash$ para & Álgebra & Topologia & Análise & Geometria \\
\midrule
Álgebra   & ---  & $16$ & $7$ & --- \\
Topologia & $25$ & ---  & --- & --- \\
Análise   & $7$  & $7$  & --- & --- \\
Geometria & $6$  & ---  & $1$ & --- \\
\bottomrule
\end{tabular}
\end{center}

\noindent Duas coisas saem daí, e nenhuma era esperada. Primeira: \textbf{Álgebra
$\leftrightarrow$ Análise é \emph{exatamente} simétrico}, $7$ numa direção e $7$ na outra --- o par
autodual não é uma figura de estilo, é uma contagem. Segunda: \emph{Topologia $\to$ Álgebra é a
aresta mais pesada} ($25$), o que confirma que a Topologia sobe da Álgebra e não o contrário.

\noindent E a Geometria cita $7$ vezes e é citada zero: \emph{é deliberado}. Ela deriva das
outras --- não elas dela ---, e uma parte que fosse citada de volta estaria a acrescentar hipóteses
em vez de as consumir. \emph{Uma folha na árvore, e não uma raiz.}

\subsection{E onde a geometria já estava, nas outras partes}\label{sec:geomapa}

\begin{center}\small
\begin{tabular}{@{}lll@{}}
\toprule
a peça & onde vive & de que lei vem \\
\midrule
a realização hiperbólica, $\ell=4\log\sigma$~\cite{series1985} & Topologia & Lei~$1$ (leitura) \\
a curva de Hilbert e as áreas (no \emph{Catálogo}) & Topologia & Lei~$1$ (leitura) \\
o reticulado de Minkowski~\cite{neukirch}     & Álgebra   & Lei~$2$ (escrita) \\
a companheira e os corpos $K_{n,m}$ (no \emph{Catálogo}) & Álgebra & Lei~$2$ (escrita) \\
a equidistribuição em $\R^{n}$ (no \emph{Catálogo})  & Análise   & Lei~$1$ (leitura) \\
\bottomrule
\end{tabular}
\end{center}

\noindent \emph{A Geometria não é uma quarta coisa a acrescentar às três}: é o que se vê quando as
duas leis se escrevem em figura. E é por isso que esta parte não abre lugar novo --- todas as suas
peças já estavam medidas noutro sítio.


\section{Conclusão das quatro partes}
\hfill \emph{\small[navegação: não afirma, situa]}
\label{sec:conc3}

As três partes deste trabalho dividem-se pelo que cada uma pode:

\begin{center}\small
\begin{tabular}{lll}
\hline
& fornece & não alcança \\ \hline
\textbf{(Parte I) álgebra} & $+$, $\times$, inverso; algoritmos exactos & a completude \\
\textbf{(Parte II) topologia} & ordem, limite, vizinhança; todo o $\R$ & as operações \\
\textbf{(Parte III) análise} & medida, taxas, distribuição & os conjuntos de medida nula \\
\hline
\end{tabular}
\end{center}

\noindent E a última linha é o remate: \emph{a análise não vê os objectos deste trabalho}. Os
corpos $K_{n,m}$, os metais, os números de Pisot --- todos de medida nula, todos invisíveis às leis
de Gauss--Kuzmin, Khinchin e Lévy, que descrevem o comportamento de \emph{quase todo} número e não
o de nenhum em particular. A análise dá as leis genéricas; a álgebra dá os casos que interessam; e
a topologia diz onde eles vivem. \emph{Nenhuma das três chega sozinha}, e essa é a razão de serem
três textos.



\begin{thebibliography}{99}
\bibitem{berger} M. Berger, \emph{Geometry Revealed: A Jacob's Ladder to Modern Higher
Geometry}, Springer, 2010.
\bibitem{coxeter} H. S. M. Coxeter, \emph{Introduction to Geometry}, 2.ª ed., Wiley, 1969.
\bibitem{cook2021} J. D. Cook, \emph{A function whose inverse is its derivative}, 2021.
\url{https://www.johndcook.com/blog/2021/09/07/when-derivative-equals-inverse/}
\bibitem{djokovic1967} D. Ž. Đoković, \emph{Product of two involutions},
Arch. Math. \textbf{18} (1967), 582--584.
\bibitem{do-carmo} M. P. do Carmo, \emph{Riemannian Geometry}, Birkhäuser, 1992.
\bibitem{amari-nagaoka} S. Amari, H. Nagaoka, \emph{Methods of Information Geometry},
Translations of Mathematical Monographs \textbf{191}, AMS, 2000.
\bibitem{cheeger-ebin} J. Cheeger, D. G. Ebin, \emph{Comparison Theorems in Riemannian
Geometry}, AMS Chelsea, 2008.
\bibitem{klein1872} F. Klein, \emph{Vergleichende Betrachtungen über neuere geometrische
Forschungen} (o Programa de Erlangen), Erlangen, 1872.
\bibitem{neukirch} J. Neukirch, \emph{Algebraic Number Theory}, Grundlehren 322, Springer, 1999.
\bibitem{rockafellar} R. T. Rockafellar, \emph{Convex Analysis}, Princeton University
Press, 1970.
\bibitem{series1985} C. Series, \emph{The modular surface and continued fractions},
J. London Math. Soc. (2) \textbf{31} (1985), 69--80.
\bibitem{myers1941} S. B. Myers, \emph{Riemannian manifolds with positive mean curvature},
Duke Math. J. \textbf{8} (1941), 401--404.
\bibitem{wonenburger1966} M. J. Wonenburger, \emph{Transformations which are products of two
involutions}, J. Math. Mech. \textbf{16} (1966), 327--338.
\bibitem{thurston} W. P. Thurston, \emph{Three-Dimensional Geometry and Topology},
Princeton University Press, 1997.
\end{thebibliography}


\end{document}
